\chapter{Complex Analysis for Computing}\label{ch:complex}

\begin{center}\itshape
The shortest path between two truths in the real domain often passes through the complex domain.\\
\upshape --- Paul Painlev\'e (1900), made famous by Jacques Hadamard
\end{center}

The integral $\int_{-\infty}^\infty\frac{\mathrm dx}{1+x^2}$ equals $\pi$. The antiderivative is $\arctan$. Now multiply the integrand by a wave: $\int_{-\infty}^\infty\frac{\cos x}{1+x^2}\dd x$. No elementary antiderivative turns up. The value is $\frac\pi e$. The number $e$ is the value of $e^{iz}$ at the point $i$, which does not lie on the real line. At that same point, $\frac1{1+z^2}$ blows up. By the end of the chapter the integral is one residue.

A complex derivative is a strong condition. The difference quotient has to settle to the same limit from every direction in the plane. A function that has such a derivative on an open set is holomorphic. It is then infinitely differentiable, it equals its Taylor series on every disc inside that set, and the values inside a circle are averages of the values on the circle. For computation, the useful consequence is smaller: the integral around a closed contour is a finite sum of residues, one number at each singular point inside.

Three promises from earlier chapters get kept here. The radius of convergence of a power series is the distance to the nearest singularity (Chapters~\ref{ch:finite} and~\ref{ch:series}). The partial fractions of $\pi\cot\pi z$, and a sixth proof that $\zeta(2)=\frac{\pi^2}6$, come from residues (Chapter~\ref{ch:series}). The Fourier transform of $\frac1{1+x^2}$ is computed again (Chapter~\ref{ch:transform}). Chapter~\ref{ch:zeta} needs one more output: $\Gamma$ as a function of a complex variable.

The contours in this chapter are boundaries of convex sets. The tools that belong to them are Goursat's lemma, Cauchy's theorem on a star-shaped set, Taylor and Laurent series, residues, and the argument principle. Cauchy's theorem in its homology form, the Phragm\'en--Lindel\"of principle, Hadamard's factorisation, and Mellin inversion are proved in Volume~8, Chapter~14. The Jordan curve theorem is Volume~7, Chapter~6. Nothing proved here depends on those five results.

\begin{think}
The function $\frac1{1+x^2}$ is infinitely differentiable on the whole real line. Its Taylor series at $0$ is $1-x^2+x^4-\cdots$, and that series converges only for $|x|<1$. What sits at distance $1$ from $0$?
\end{think}

\section{Holomorphic functions and power series}

Throughout, $U$ is an open subset of $\C$. Write $D(a,r)=\{z:|z-a|<r\}$ for the open disc and $\overline D(a,r)$ for the closed disc. We identify $\C$ with $\R^2$, so continuity, open sets, and compactness keep their usual meaning.

\begin{definition}
A function $f:U\to\C$ is \emph{complex differentiable} at $z_0\in U$ if
\[
f'(z_0)=\lim_{h\to0}\frac{f(z_0+h)-f(z_0)}h
\]
exists, with $h$ tending to $0$ in $\C$. It is \emph{holomorphic} on $U$ if it is complex differentiable at every point of $U$. A function holomorphic on all of $\C$ is \emph{entire}.
\end{definition}

Sums, products, quotients, and compositions follow the same rules as for real functions, with the same proofs. Polynomials are entire. A rational function is holomorphic off the zeros of its denominator. The new demand is directional: $h$ may approach $0$ along any path, and the limit has to be the same limit.

\begin{proposition}[Cauchy--Riemann equations]\label{prop:cr}
Write $f=u+iv$ with $u,v$ real, as functions of $(x,y)$ with $z=x+iy$. If $f$ is complex differentiable at $z_0$, then the partial
derivatives of $u$ and $v$ exist at $z_0$ and
\[
u_x=v_y,\qquad u_y=-v_x,\qquad f'(z_0)=u_x+iv_x.
\]
\end{proposition}
\begin{proof}
Let $h$ tend to $0$ along the real axis. The difference quotient tends to $u_x+iv_x$. Along the imaginary axis, take $h=ik$ with $k$ real. The quotient tends to $\frac{u_y+iv_y}i=v_y-iu_y$. Both limits equal $f'(z_0)$. Match real and imaginary parts.
\end{proof}

So $\overline z$ fails the test at every point: $u=x$, $v=-y$, and $u_x=1$ while $v_y=-1$. The function $|z|^2$ passes only at $0$. Both are smooth as functions of two real variables. The converse is true as well: if $u$ and $v$ have continuous partial derivatives and those derivatives satisfy the Cauchy--Riemann equations, then $f$ is holomorphic. This chapter does not call on that converse.

\subsection*{Power series}

\begin{lemma}[Abel]\label{lem:abel}
Let $(c_n)$ be complex numbers and set
\[
R=\sup\{r\ge0:(|c_n|r^n)_n\text{ is bounded}\}\in[0,\infty].
\]
Then $\sum c_nz^n$ converges
absolutely for $|z|<R$, uniformly on every closed disc $\overline D(0,r)$ with $r<R$, and its terms are unbounded for $|z|>R$.
\end{lemma}
\begin{proof}
\textit{Outside the disc.} If $|z|>R$, the terms $|c_nz^n|$ are unbounded, by the definition of $R$. A series with unbounded terms diverges.

\textit{Inside a smaller disc.} Fix $r<\rho<R$. Some $M$ satisfies $|c_n|\rho^n\le M$ for every $n$. On the closed disc $|z|\le r$ the general term is then bounded by
\[
|c_nz^n|\le M\Bigl(\frac r\rho\Bigr)^n.
\]
The right-hand side is a convergent geometric series, and it does not depend on $z$. The Weierstrass test (Lemma~\ref{lem:uniform}; the same proof works on any set) gives absolute and uniform convergence on $\overline D(0,r)$.
\end{proof}

The number $R$ is the \emph{radius of convergence}. Inside $D(0,R)$ the series converges. Outside, the terms do not tend to $0$, so the series diverges. On the circle $|z|=R$, anything can happen.

\begin{theorem}\label{thm:powerseries}
Let $f(z)=\sum_{n\ge0}c_n(z-a)^n$ have radius of convergence $R>0$. Then $f$ is holomorphic on $D(a,R)$ and
\[
f'(z)=\sum_{n\ge1}nc_n(z-a)^{n-1},
\]
a series with the same radius. Consequently $f$ has derivatives of all orders on $D(a,R)$, and $c_n=\frac{f^{(n)}(a)}{n!}$.
\end{theorem}
\begin{proof}
Take $a=0$. \textit{The derived series.} If $r<\rho<R$ and $|c_n|\rho^n\le M$, then
\[
n|c_n|r^{n-1}\le\frac M\rho\,n\Bigl(\frac r\rho\Bigr)^{n-1}.
\]
The series $\sum n t^{n-1}$ converges for $0\le t<1$, so the right-hand side is summable. The derived series converges absolutely on $\overline D(0,r)$. Its radius is at least $R$. It is at most $R$, because $|c_nz^n|\le|z|\cdot|nc_nz^{n-1}|$ for $n\ge1$.

\textit{The difference quotient.} Fix $z$ with $|z|<r<R$. For $|w|\le r$ set $q_n(w)=\sum_{k=0}^{n-1}w^kz^{n-1-k}$. Then $q_n(w)(w-z)=w^n-z^n$, $q_n(z)=nz^{n-1}$, and $|q_n(w)|\le nr^{n-1}$. The series $Q(w)=\sum_nc_nq_n(w)$ converges uniformly on $\overline D(0,r)$, by the bound of the previous paragraph, so $Q$ is continuous there. For $w\ne z$,
\[
Q(w)=\frac{f(w)-f(z)}{w-z}.
\]
Let $w\to z$. Continuity of $Q$ gives $f'(z)=Q(z)=\sum nc_nz^{n-1}$.

\textit{Higher derivatives.} Apply the same result to $f'$, then to $f''$, and so on:
\[
f^{(k)}(z)=\sum_{n\ge k}n(n-1)\cdots(n-k+1)c_nz^{n-k}.
\]
At $z=0$ this is $f^{(k)}(0)=k!\,c_k$.
\end{proof}

\begin{example}[exponential and logarithm]
The series $\exp z=\sum\frac{z^n}{n!}$ has infinite radius, so $\exp$ is entire and $\exp'=\exp$. The identity $e^{a+b}=e^ae^b$ (Cauchy product, Volume~1, Chapter~3) and Euler's formula (Volume~1, Chapter~5) give $e^{x+iy}=e^x(\cos y+i\sin y)$. So $\exp$ never vanishes, and $e^z=1$ exactly when $z\in2\pi i\Z$. The same formulas give $\cos z=\frac{e^{iz}+e^{-iz}}2$ and $\sin z=\frac{e^{iz}-e^{-iz}}{2i}$, both entire. And $\sin z=0$ means $e^{2iz}=1$, that is $z\in\pi\Z$. The sine has no zeros off the real axis, so $\cot z=\frac{\cos z}{\sin z}$ is holomorphic on $\C\setminus\pi\Z$.

On the \emph{slit plane} $\C\setminus\left]-\infty,0\right]$, every $z$ is $re^{i\theta}$ with $r>0$ and a unique $\theta\in\left]-\pi,\pi\right[$. The half-angle identity $\frac{\sin\theta}{1+\cos\theta}=\tan\frac\theta2$ gives the continuous formula $\theta=2\arctan\frac{\operatorname{Im}z}{|z|+\operatorname{Re}z}$. The \emph{principal logarithm} $\operatorname{Log}z=\log r+i\theta$ is therefore continuous, and it inverts $\exp$. It is holomorphic, with derivative $\frac1z$. Write $\ell=\operatorname{Log}z$ and $\ell_w=\operatorname{Log}w$. Then $\frac{\ell_w-\ell}{w-z}=\frac{\ell_w-\ell}{e^{\ell_w}-e^\ell}\to\frac1{e^\ell}=\frac1z$ as $w\to z$, because $\ell_w\to\ell$. Powers are $z^s=e^{s\operatorname{Log}z}$ on the slit plane.
\end{example}

Check the exponential at one point. At $z=i\pi$, Euler's formula gives $e^{i\pi}=\cos\pi+i\sin\pi=-1$. The series $\sum (i\pi)^n/n!$ is an entire function, so it has to take that same value.

\section{Path integrals and Cauchy's theorem}

A \emph{path} is a continuous, piecewise $C^1$ map $\gamma:[\alpha,\beta]\to\C$. It is \emph{closed} if $\gamma(\alpha)=\gamma(\beta)$. Write $\gamma^*$ for its image. If $f$ is continuous on $\gamma^*$,
\[
\begin{gathered}\int_\gamma f(z)\dd z=\int_\alpha^\beta f(\gamma(t))\gamma'(t)\dd t,\qq L(\gamma)=\int_\alpha^\beta|\gamma'(t)|\dd t.\end{gathered}
\]
An increasing piecewise $C^1$ change of parameter leaves the integral unchanged, by substitution. Reversing the path changes the sign. Concatenation adds the integrals. We write $[z,w]$ for the segment from $z$ to $w$, and $C(a,r)$ for the circle $t\mapsto a+re^{it}$, $t\in[0,2\pi]$. The symbol $\partial T$ means the boundary of a triangle $T$, traversed once along each side.

\begin{proposition}\label{prop:pathint}
\begin{enumerate}[label=(\alph*)]
\item $\bigl|\int_\gamma f\bigr|\le L(\gamma)\max_{\gamma^*}|f|$. Hence if $f_n\to f$ uniformly on $\gamma^*$, then $\int_\gamma f_n\to\int_\gamma f$.
\item If $F$ is holomorphic on an open set containing $\gamma^*$ and $F'=f$ is continuous, then $\int_\gamma f=F(\gamma(\beta))-F(\gamma(\alpha))$. In
particular $\int_\gamma f=0$ for every closed path.
\item For $n\in\Z$, $\displaystyle\int_{C(a,r)}(z-a)^n\dd z=0$ if $n\ne-1$, and $\displaystyle\int_{C(a,r)}\frac{\mathrm dz}{z-a}=2\pi i$.
\end{enumerate}
\end{proposition}
\begin{proof}
(a) For a complex-valued integrand $g$, $|\int_\alpha^\beta g|\le\int_\alpha^\beta|g|$. Multiply $g$ by a unit complex number that makes the integral real, and use the real inequality. The length bound is that estimate with $g=f(\gamma)\gamma'$. Uniform convergence passes inside a fixed integral because the difference is at most $L(\gamma)$ times the uniform distance.

(b) On each interval where $\gamma$ is $C^1$, $\frac{\mathrm d}{\mathrm dt}F(\gamma(t))=F'(\gamma(t))\gamma'(t)$, by the same proof as the real chain rule. Integrate on each piece and add.

(c) For $n\ne-1$, $\frac{(z-a)^{n+1}}{n+1}$ is a primitive on $\C\setminus\{a\}$, so (b) gives $0$ on the closed circle. For $n=-1$,
\[
\int_0^{2\pi}\frac{ire^{it}}{re^{it}}\dd t=\int_0^{2\pi}i\dd t=2\pi i.
\]
\end{proof}

The function $\frac1z$ has a primitive on the slit plane, namely $\operatorname{Log}z$. It has none on $\C\setminus\{0\}$: the integral around the unit circle is $2\pi i$, not $0$. Go around the origin and the logarithm does not come back to its starting value. Every theorem below grows from that one integral.

\begin{theorem}[Goursat]\label{thm:goursat}
Let $f$ be holomorphic on $U$ and $T\subseteq U$ a closed triangle. Then $\int_{\partial T}f=0$.
\end{theorem}
\begin{proof}
\textit{Setup.} Let $I=\int_{\partial T}f$. Let $d$ be the diameter of $T$ and $p$ its perimeter.

\textit{Subdivision.} Join the midpoints of the sides. That cuts $T$ into four triangles, each similar to $T$ with ratio $\frac12$ (Figure~\ref{fig:goursat}). Orient each small boundary like $T$. The four boundary integrals add up to $I$: each inner side is walked once in each direction, so those contributions cancel. One of the four triangles, call it $T_1$, carries at least a quarter of the integral:
\[
\Bigl|\int_{\partial T_1}f\Bigr|\ge\frac{|I|}4.
\]
Repeat. You get nested triangles $T\supseteq T_1\supseteq T_2\supseteq\cdots$ with diameter $\frac d{2^k}$, perimeter $\frac p{2^k}$, and
\[
|I|\le 4^k\Bigl|\int_{\partial T_k}f\Bigr|.
\]

\textit{Estimate.} The sets $T_k$ are nested, non-empty, and compact, so they share a point $z_0\in T\subseteq U$. Since $f$ is complex differentiable at $z_0$,
\[
f(z)=f(z_0)+f'(z_0)(z-z_0)+\varepsilon(z)(z-z_0),\qquad \varepsilon(z)\to0\text{ as }z\to z_0.
\]
The polynomial part $f(z_0)+f'(z_0)(z-z_0)$ has a primitive, so its integral over $\partial T_k$ is $0$ by Proposition~\ref{prop:pathint}(b). What remains is the integral of $\varepsilon(z)(z-z_0)$. On $T_k$ one has $|z-z_0|\le\frac d{2^k}$, and the path $\partial T_k$ has length $\frac p{2^k}$. Proposition~\ref{prop:pathint}(a) bounds that integral by
\[
\frac p{2^k}\cdot\frac d{2^k}\cdot\sup_{T_k}|\varepsilon|.
\]
Multiply by $4^k$:
\[
|I|\le pd\sup_{T_k}|\varepsilon|.
\]
The quantity being bounded is $|I|$, the integral over the original triangle. The product $pd$ is fixed. As $k\to\infty$ the triangles shrink to $z_0$, so $\sup_{T_k}|\varepsilon|\to0$.

\textit{Finish.} The bound tends to $0$. Therefore $I=0$.
\end{proof}

\begin{figure}[ht]
\centering
\fig{ch09-goursat}
\caption{Goursat's proof. Of the four midpoint triangles, keep one carrying at least a quarter of the integral, and repeat. The
triangles shrink to a point $z_0$ where $f$ is almost linear, and linear functions have integral $0$.}
\label{fig:goursat}
\end{figure}

Goursat's note of 1900 assumes nothing about $f'$, not even continuity. Cauchy's earlier proofs had asked for a continuous derivative. The subdivision replaces that hypothesis: once the triangles are small, holomorphy at the single point $z_0$ makes $f$ close to a linear function, and a linear function integrates to zero around a closed path.

A version with one bad point is needed for residues.

\begin{corollary}\label{cor:goursatpoint}
The conclusion of Theorem~\ref{thm:goursat} holds if $f$ is continuous on $U$ and holomorphic on $U\setminus\{p\}$ for one point $p$.
\end{corollary}
\begin{proof}
If $p\notin T$, this is the theorem. If $p$ is a vertex of $T$, choose points $q,q'$ on the two sides issuing from $p$, at distance at most $\varepsilon$ from $p$. The triangle $T$ splits into the small triangle $pqq'$ and two triangles that do not contain $p$. The integrals over those two vanish. The integral over $\partial T$ is therefore the integral over $\partial(pqq')$. Its modulus is at most the perimeter of $pqq'$, which is at most $4\varepsilon$, times $\max_T|f|$. The quantity being bounded is $\bigl|\int_{\partial T}f\bigr|$. The bound $4\varepsilon\max_T|f|$ tends to $0$ as $\varepsilon\to0$, because $f$ is continuous on the compact triangle, hence bounded. If $p$ is any other point of $T$, cut $T$ into two or three triangles that have $p$ as a vertex, and apply the vertex case to each.
\end{proof}

An open set $U$ is \emph{star-shaped} with respect to a point $c\in U$ when the segment $[c,z]$ lies in $U$ for every $z\in U$. A convex set is star-shaped with respect to each of its points. Discs and half-planes are convex. The slit plane is star-shaped with respect to $c=1$: the ray $(-\infty,0]$ is the only cut, and it does not meet any segment from $1$ to a point off that ray.

Figure~\ref{fig:rw09-star} shows a set that is star-shaped and not convex. The five lobes bend in, so some chords between boundary points leave the set. Every segment that starts at $c$ stays inside. One triangle with a vertex at $c$ is cut into the four midpoint triangles of Goursat's argument. Each of those small triangles lies in the set, so the integral of a holomorphic function over its boundary is zero.

\begin{figure}[ht]
\centering
\fig{ch09-rewrite-star}
\caption{A set star-shaped from $c$: every ray from $c$ stays inside. One triangle $c$, $z$, $z{+}h$ is split into the four midpoint triangles used by Goursat.}
\label{fig:rw09-star}
\end{figure}

Augustin-Louis Cauchy's memoir of 1825 evaluated real integrals by paths that leave the real line. The form proved here is the one Volume~8 cites: on a star-shaped set, a continuous function that is holomorphic except at one point has a primitive $F$ with $F'=f$, and the integral over every closed path in the set is zero.

\portrait{cauchy}{Augustin-Louis Cauchy (1789--1857). Illustration from \emph{Das neunzehnte Jahrhundert in Bildnissen}, Berlin, 1901. Public domain.}

\begin{theorem}[Cauchy's theorem]\label{thm:cauchy}
Let $U$ be star-shaped with respect to $c$, and $f$ continuous on $U$ and holomorphic on $U\setminus\{p\}$. Then $F(z)=\int_{[c,z]}f$ is
holomorphic on $U$ with $F'=f$, and $\int_\gamma f=0$ for every closed path $\gamma$ in $U$.
\end{theorem}
\begin{proof}
\textit{Setup.} Fix $z\in U$ and a radius $r>0$ with $D(z,r)\subseteq U$. For $|h|<r$ the segment $[z,z+h]$ lies in that disc. Every point of the segment is joined to $c$ by a segment that lies in $U$, because $U$ is star-shaped and the disc sits inside $U$. The triangle with vertices $c$, $z$, $z+h$ is therefore contained in $U$.

\textit{The difference quotient.} Corollary~\ref{cor:goursatpoint} gives $\int_{\partial T}f=0$ on that triangle, so
\[
F(z+h)-F(z)=\int_{[z,z+h]}f.
\]
The error to estimate is
\[
\Delta=\int_{[z,z+h]}\bigl(f(w)-f(z)\bigr)\dd w,
\]
which is $F(z+h)-F(z)-hf(z)$. Proposition~\ref{prop:pathint}(a) bounds the difference quotient:
\[
\begin{aligned}{}&\Bigl|\frac{F(z+h)-F(z)}h-f(z)\Bigr|\eqbr\le\max_{w\in[z,z+h]}|f(w)-f(z)|.\end{aligned}
\]
The quantity being bounded is the difference between the difference quotient of $F$ and the value $f(z)$. The right-hand side tends to $0$ as $h\to0$, because $f$ is continuous at $z$ and every $w$ on the segment is within $|h|$ of $z$. Thus $F'(z)=f(z)$.

\textit{Closed paths.} The derivative $F'=f$ is continuous, so Proposition~\ref{prop:pathint}(b) gives $\int_\gamma f=0$ for every closed path in $U$.
\end{proof}

\begin{remark}
The same vanishing holds on sets that are not star-shaped, as soon as the closed path winds around no point outside the domain. That homology form of Cauchy's theorem is proved in Volume~8, Chapter~14, together with the Phragm\'en--Lindel\"of principle, Hadamard's factorisation, and Mellin inversion. The Jordan curve theorem, which gives a simple closed curve an inside and an outside, is Volume~7, Chapter~6. Every contour used below is the boundary of a convex set, and Theorem~\ref{thm:cauchy} covers those.
\end{remark}

\section{Cauchy's formula and its consequences}

\begin{theorem}[Cauchy's integral formula]\label{thm:cauchyformula}
Let $f$ be holomorphic on $U$ and $\overline D(a,r)\subseteq U$. Then for $|z-a|<r$,
\[
f(z)=\frac1{2\pi i}\int_{C(a,r)}\frac{f(w)}{w-z}\dd w.
\]
\end{theorem}
\begin{proof}
\textit{The kernel alone.} First, $\int_{C(a,r)}\frac{\mathrm dw}{w-z}=2\pi i$. On the circle,
\[
\frac1{w-z}=\frac1{(w-a)-(z-a)}=\sum_{n\ge0}\frac{(z-a)^n}{(w-a)^{n+1}}.
\]
The general term is bounded by $\frac1r\bigl(\frac{|z-a|}r\bigr)^n$, a convergent geometric series independent of $w$ on the circle. Convergence is uniform. Integrate term by term. Proposition~\ref{prop:pathint}(c) says that only the term $n=0$ survives, and it contributes $2\pi i$.

\textit{Remove the value $f(z)$.} The closed disc is compact and $\C\setminus U$ is closed, so the distance between them is positive. Some $r'>r$ has $D(a,r')\subseteq U$. Define $g(w)=\frac{f(w)-f(z)}{w-z}$ for $w\ne z$, and $g(z)=f'(z)$. Then $g$ is continuous on $D(a,r')$ and holomorphic except at the one point $z$. Theorem~\ref{thm:cauchy}, applied on a disc, which is star-shaped, gives $\int_{C(a,r)}g=0$. That is
\[
\int_{C(a,r)}\frac{f(w)}{w-z}\dd w=f(z)\int_{C(a,r)}\frac{\mathrm dw}{w-z}=2\pi if(z).
\]
\end{proof}

The values inside the disc are averages of the values on the circle. At the centre, the formula collapses to
\[
f(a)=\frac1{2\pi}\int_0^{2\pi}f(a+re^{it})\dd t.
\]
This is the \emph{mean value property}. The kernel $\frac1{w-z}$ is a geometric series in the variable $z$, so the same integral produces the Taylor series.

\begin{theorem}[holomorphic means analytic]\label{thm:analytic}
Let $f$ be holomorphic on $U$ and $D(a,R)\subseteq U$. Then $f$ is the sum of its Taylor series on $D(a,R)$:
\[
\begin{gathered}f(z)=\sum_{n\ge0}c_n(z-a)^n,\qq c_n=\frac{f^{(n)}(a)}{n!}=\frac1{2\pi i}\int_{C(a,r)}\frac{f(w)\dd w}{(w-a)^{n+1}}\end{gathered}
\]
for every $r\in\left]0,R\right[$.
In particular $f$ has derivatives of all orders, all holomorphic on $U$.
\end{theorem}
\begin{proof}
Fix $0<r<R$ and $|z-a|<r$. Expand $\frac1{w-z}$ as in the previous proof, uniformly for $w$ on $C(a,r)$. Integrate term by term in Cauchy's formula. You get $f(z)=\sum c_n(r)(z-a)^n$, where $c_n(r)$ is the integral in the statement. So $f$ is given by a power series on $D(a,r)$. Theorem~\ref{thm:powerseries} gives $c_n(r)=\frac{f^{(n)}(a)}{n!}$, which does not depend on $r$. Every point of $D(a,R)$ lies in some $D(a,r)$ with $r<R$. Near each point of $U$, the derivative $f'$ is itself a power series, hence holomorphic.
\end{proof}

\begin{corollary}[Cauchy's estimates]\label{cor:radius}
\begin{enumerate}[label=(\alph*)]
\item With $M_r=\max_{C(a,r)}|f|$, $\ \bigl|f^{(n)}(a)\bigr|\le\frac{n!\,M_r}{r^n}$.
\item The Taylor series of $f$ at $a$ has radius of convergence at least the radius of the largest open disc around $a$ contained in
$U$. If $f$ is unbounded near a point $b$ of the boundary of that disc, the radius is exactly $|b-a|$.
\end{enumerate}
\end{corollary}
\begin{proof}
(a) The coefficient integral has length $2\pi r$ and integrand bounded by $M_r/r^{n+1}$. Proposition~\ref{prop:pathint}(a) gives
\[
|c_n|\le\frac1{2\pi}\cdot 2\pi r\cdot\frac{M_r}{r^{n+1}}=\frac{M_r}{r^n},
\]
and $|f^{(n)}(a)|=n!\,|c_n|$. The quantity being bounded is the derivative at the single point $a$.

(b) The lower bound on the radius is Theorem~\ref{thm:analytic}. If the radius were larger than $|b-a|$, the sum of the series would be continuous at $b$, hence bounded near $b$, and on $D(a,|b-a|)$ that sum equals $f$. Then $f$ would be bounded near $b$ inside the disc.
\end{proof}

\begin{example}
This answers the pause box. The function $\frac1{1+z^2}$ is holomorphic on $\C$ except at $\pm i$, and it is unbounded at each of those two points. The largest open disc about $0$ that avoids them has radius $1$. Corollary~\ref{cor:radius}(b) says the Taylor series at $0$ has radius exactly $1$. The obstruction sits at $i$ and at $-i$. The real line does not pass through either point, so a real-variable test never sees it.

The same reading gives the radius for the Bernoulli generating function. Write $\frac z{e^z-1}=\frac1{h(z)}$ with $h(z)=\sum_{n\ge0}\frac{z^n}{(n+1)!}$. The function $h$ is entire. It vanishes where $e^z=1$ and $z\ne0$, that is on $2\pi i\Z\setminus\{0\}$. So $\frac z{e^z-1}$ is holomorphic on $D(0,2\pi)$ and unbounded near $\pm2\pi i$. The radius is exactly $2\pi$, as Chapters~\ref{ch:finite} and~\ref{ch:series} announced.
\end{example}

A continuous function whose integrals over triangles all vanish is holomorphic. That converse of Goursat is what makes holomorphy stable under limits.

\begin{theorem}[Morera]\label{thm:morera}
Let $f$ be continuous on $U$ with $\int_{\partial T}f=0$ for every closed triangle $T\subseteq U$. Then $f$ is holomorphic on $U$.
\end{theorem}
\begin{proof}
Take a disc $D\subseteq U$ centred at $c$. The proof of Theorem~\ref{thm:cauchy} used only the vanishing of triangle integrals, which we have by hypothesis. It produces $F(z)=\int_{[c,z]}f$, holomorphic on $D$, with $F'=f$. Theorem~\ref{thm:analytic} says that $F'$ is holomorphic on $D$.
\end{proof}

\begin{theorem}[limits and integrals]\label{thm:holoparam}
\begin{enumerate}[label=(\alph*)]
\item (Weierstrass) If $f_n$ are holomorphic on $U$ and converge to $f$ uniformly on every closed disc contained in $U$, then $f$ is
holomorphic, and $f_n'\to f'$ uniformly on every closed disc contained in $U$. In particular a series of holomorphic functions converging
uniformly on such discs may be differentiated term by term.
\item Let $I$ be an interval and $f$ continuous on $I\times U$, holomorphic in $z$ for each $t\in I$. Suppose that for every closed disc
$K\subseteq U$ there is a continuous integrable $g_K$ on $I$ with $|f(t,z)|\le g_K(t)$ for $z\in K$. Then $F(z)=\int_If(t,z)\dd t$ is holomorphic on $U$.
\end{enumerate}
\end{theorem}
\begin{proof}
(a) \textit{The limit is holomorphic.} A locally uniform limit of continuous functions is continuous, so $f$ is continuous. Let $T$ be a triangle in a disc $D\subseteq U$. Proposition~\ref{prop:pathint}(a) lets the uniform limit pass inside the integral, and Goursat gives $\int_{\partial T}f_n=0$, so $\int_{\partial T}f=0$. Morera says $f$ is holomorphic.

\textit{The derivatives converge.} Let $\overline D(a,r)\subseteq U$, and choose $\rho>0$ with $\overline D(a,r+2\rho)\subseteq U$. For $|z-a|\le r$, Theorem~\ref{thm:analytic} on the disc $D(z,\rho)$ gives
\[
f_n'(z)-f'(z)=\frac1{2\pi i}\int_{C(z,\rho)}\frac{f_n(w)-f(w)}{(w-z)^2}\dd w.
\]
The path has length $2\pi\rho$, and $|w-z|=\rho$ on it, so the modulus of the integral is at most
\[
\frac1\rho\max_{\overline D(a,r+\rho)}|f_n-f|.
\]
The quantity being bounded is $|f_n'(z)-f'(z)|$, uniformly for $|z-a|\le r$. The maximum of $|f_n-f|$ on that closed disc tends to $0$ by hypothesis. So $f_n'\to f'$ uniformly on $\overline D(a,r)$.

(b) \textit{Continuity of $F$.} Let $z_n\to z$ inside a closed disc $K\subseteq U$. On every closed bounded subinterval of $I$, the integrand $f(t,z_n)$ tends to $f(t,z)$ uniformly in $t$, by uniform continuity on a compact set. The domination $|f|\le g_K$ and Corollary~\ref{cor:dct} give $F(z_n)\to F(z)$.

\textit{Triangle integrals.} Let $T$ be a triangle in a disc contained in $U$, with $\partial T$ made of three segments parametrised on $[0,1]$, and let $K$ be a closed disc containing $T$. For a closed bounded interval $[u,v]\subseteq I$, Fubini on the rectangle $[u,v]\times[0,1]$ (Theorem~\ref{thm:fubinirect}) gives
\[
\int_{\partial T}\int_u^vf(t,z)\dd t\dd z=\int_u^v\int_{\partial T}f(t,z)\dd z\dd t=0,
\]
the inner integral vanishing by Goursat. As $[u,v]$ grows to cover $I$, the left side tends to $\int_{\partial T}F$. The difference is at most
\[
L(\partial T)\int_{I\setminus[u,v]}g_K.
\]
The quantity being bounded is that difference. The length $L(\partial T)$ is fixed, and the integral of $g_K$ over the complement tends to $0$ because $g_K$ is integrable. So $\int_{\partial T}F=0$. Morera says $F$ is holomorphic.
\end{proof}

\subsection*{Liouville, zeros and the maximum principle}

\begin{theorem}[Liouville]\label{thm:liouville}
A bounded entire function is constant.
\end{theorem}
\begin{proof}
Suppose $|f|\le M$ on $\C$. Theorem~\ref{thm:analytic} writes $f(z)=\sum c_nz^n$ on $\C$. Corollary~\ref{cor:radius}(a), on the circle of radius $r$, gives $|c_n|\le M/r^n$ for every $r>0$. Fix $n\ge1$. The quantity being bounded is the single coefficient $c_n$. The numerator $M$ stays fixed while the denominator $r^n$ tends to infinity as $r\to\infty$. The bound tends to $0$, so $c_n=0$.
\end{proof}

\begin{corollary}[fundamental theorem of algebra]\label{cor:fta}
Every non-constant complex polynomial has a root in $\C$.
\end{corollary}
\begin{proof}
Let $p(z)=a_nz^n+\dots+a_0$ with $n\ge1$ and $a_n\ne0$. Then
\[
|p(z)|\ge|a_n||z|^n\Bigl(1-\sum_{k<n}\frac{|a_k|}{|a_n||z|^{n-k}}\Bigr)\to\infty
\]
as $|z|\to\infty$. If $p$ had no root, $\frac1p$ would be entire. It tends to $0$ at infinity, so it is bounded: continuous on a large closed disc, and small outside that disc. Liouville would make $\frac1p$ constant, and then $p$ would be constant. The degree is at least $1$, so a root exists.
\end{proof}

Volume~1, Chapter~6, proved the same theorem by watching the minimum of $|p|$, which is d'Alembert's argument. Liouville gives it in the two lines above. Section~\ref{sec:argument} gives a third proof, and that one also counts the roots.

An open set $U$ is \emph{connected} when any two of its points can be joined by a path in $U$.

\begin{theorem}[zeros and the identity theorem]\label{thm:identity}
Let $f$ be holomorphic on a connected open set $U$, and not identically zero.
\begin{enumerate}[label=(\alph*)]
\item For every $a\in U$ there are $m\ge0$ and $g$ holomorphic on $U$ with $g(a)\ne0$ and $f(z)=(z-a)^mg(z)$. In particular the zeros of $f$ are isolated.
\item Consequently, two holomorphic functions on $U$ that agree on a set having an accumulation point in $U$ are equal.
\end{enumerate}
\end{theorem}
\begin{proof}
Let $A$ be the set of points $z\in U$ at which every derivative $f^{(n)}(z)$, $n\ge0$, vanishes. The set $A$ is open: near a point of $A$, Theorem~\ref{thm:analytic} says $f$ equals its Taylor series, and that series is $0$, so $f$ and all its derivatives vanish on a disc. The complement $U\setminus A$ is open as well: if $f^{(n)}(z)\ne0$, continuity keeps $f^{(n)}\ne0$ nearby. Suppose both $A$ and $U\setminus A$ are non-empty. A path $\gamma:[0,1]\to U$ from a point of $A$ to a point of $U\setminus A$ would split $[0,1]$ into the two non-empty disjoint relatively open sets $\gamma^{-1}(A)$ and $\gamma^{-1}(U\setminus A)$. That split is impossible: $[0,1]$ is connected. (If $0$ lies in the first set, the supremum of that set is a point that cannot lie in either piece.) Since $f$ is not identically zero, $A$ is empty.

(a) At $a$, let $m$ be the least $n$ with $f^{(n)}(a)\ne0$. Near $a$,
\[
f(z)=\sum_{n\ge m}c_n(z-a)^n=(z-a)^m\sum_{k\ge0}c_{m+k}(z-a)^k.
\]
Set $g(z)=\frac{f(z)}{(z-a)^m}$ for $z\ne a$, and $g(a)=c_m\ne0$. Near $a$, $g$ is a power series, so $g$ is holomorphic, and $g(a)\ne0$ keeps $g$ non-zero on a small disc. That disc contains no other zero of $f$.

(b) Apply (a) to the difference of the two functions. An accumulation point of zeros is not an isolated zero, so the difference is identically zero.
\end{proof}

An identity between holomorphic functions, once proved on a real interval, holds on any connected open set where both sides are holomorphic. For instance $\sin^2z+\cos^2z=1$ on all of $\C$. Section~\ref{sec:gammac} uses the same transfer for $\Gamma$.

\begin{theorem}[maximum principle]\label{thm:maxprinciple}
If $f$ is holomorphic on a connected open set $U$ and $|f|$ has a local maximum at $a\in U$, then $f$ is constant.
\end{theorem}
\begin{proof}
Suppose $|f|\le|f(a)|$ on $D(a,R)$, and write $f(a)=|f(a)|e^{i\alpha}$. For $0<r<R$ the mean value property gives
\[
|f(a)|=\frac1{2\pi}\int_0^{2\pi}\operatorname{Re}\bigl(e^{-i\alpha}f(a+re^{it})\bigr)\dd t.
\]
The integrand is at most $|f(a+re^{it})|\le|f(a)|$. A continuous function that stays at most a constant, and whose average equals that constant, equals the constant everywhere. So $\operatorname{Re}(e^{-i\alpha}f)=|f(a)|$ on the circle. Combined with $|f|\le|f(a)|$, this forces $e^{-i\alpha}f=|f(a)|$ on the circle. Thus $f=f(a)$ on $D(a,R)$. Theorem~\ref{thm:identity} spreads the equality to all of $U$.
\end{proof}

On a bounded open set, if $|f|$ is continuous up to the boundary, the maximum of $|f|$ is attained on the boundary. The averaging is the same one as the maximum principle in miniature of Chapter~\ref{ch:series}.
\section{Singularities, Laurent series and residues}

Let $f$ be holomorphic on a punctured disc $D(a,r)\setminus\{a\}$. The point $a$ is an \emph{isolated singularity}.

\begin{theorem}[Riemann's removable singularity theorem]\label{thm:removable}
If $f$ is bounded near $a$, it extends to a holomorphic function on $D(a,r)$.
\end{theorem}
\begin{proof}
Set $h(z)=(z-a)^2f(z)$ for $z\ne a$, and $h(a)=0$. Then
\[
\frac{h(z)-h(a)}{z-a}=(z-a)f(z).
\]
The quantity $(z-a)f(z)$ tends to $0$ as $z\to a$, because $f$ is bounded near $a$ and the factor $z-a$ tends to $0$. So $h$ is complex differentiable at $a$, with $h'(a)=0$, and holomorphic on the rest of the disc. Theorem~\ref{thm:analytic} writes $h(z)=\sum_{n\ge2}c_n(z-a)^n$, since $h(a)=h'(a)=0$. The series $\sum_{n\ge2}c_n(z-a)^{n-2}$ is the holomorphic extension of $f$.
\end{proof}

\begin{lemma}\label{lem:annulus}
Let $g$ be holomorphic on the annulus $r<|z-a|<R$. Then $\int_{C(a,\rho)}g$ does not depend on $\rho\in\left]r,R\right[$.
\end{lemma}
\begin{proof}
Write $\phi(\rho,t)=g(a+\rho e^{it})\,i\rho e^{it}$, so the integral is $J(\rho)=\int_0^{2\pi}\phi(\rho,t)\dd t$. Theorem~\ref{thm:analytic} says $g'$ is continuous. The chain rule gives
\[
\begin{aligned}\partial_\rho\phi&=g'(a+\rho e^{it})\,i\rho e^{2it}+g(a+\rho e^{it})\,ie^{it}\eqbr=\partial_t\bigl(g(a+\rho e^{it})e^{it}\bigr).\end{aligned}
\]
Theorem~\ref{thm:paramcompact} lets the derivative pass under the integral:
\[
J'(\rho)=\int_0^{2\pi}\partial_t\bigl(g(a+\rho e^{it})e^{it}\bigr)\dd t.
\]
The integrand is a derivative in $t$ of a function that is $2\pi$-periodic, so the integral is zero. The quantity $J'(\rho)$ is the rate of change of the circle integral as the radius moves. It vanishes, so $J(\rho)$ is constant on $\left]r,R\right[$.
\end{proof}

\begin{theorem}[Laurent]\label{thm:laurent}
Let $f$ be holomorphic on the annulus $A=\{r<|z-a|<R\}$, with $0\le r<R\le\infty$. Then
\[
\begin{gathered}f(z)=\sum_{n\in\Z}c_n(z-a)^n,\qq c_n=\frac1{2\pi i}\int_{C(a,\rho)}\frac{f(w)\dd w}{(w-a)^{n+1}}\quad(r<\rho<R),\end{gathered}
\]
where the coefficients do not depend on $\rho$. The part $\sum_{n\ge0}$ converges on $D(a,R)$; the part $\sum_{n<0}$ converges for $|z-a|>r$ and is
holomorphic there; both converge uniformly on every closed annulus $\rho_1\le|z-a|\le\rho_2$ inside $A$. The coefficients are unique.
\end{theorem}
\begin{proof}
\textit{The coefficients.} Independence of $\rho$ is Lemma~\ref{lem:annulus} applied to $\frac{f(w)}{(w-a)^{n+1}}$.

\textit{Two circles.} Fix $z$ and radii $r<\rho_1<|z-a|<\rho_2<R$. The function $g(w)=\frac{f(w)-f(z)}{w-z}$ is bounded near $z$, so Theorem~\ref{thm:removable} makes it holomorphic on the annulus. Lemma~\ref{lem:annulus} gives $\int_{C(a,\rho_2)}g=\int_{C(a,\rho_1)}g$. On the outer circle, the proof of Theorem~\ref{thm:cauchyformula} gives $\int_{C(a,\rho_2)}\frac{\mathrm dw}{w-z}=2\pi i$. On the inner circle,
\[
\frac1{w-z}=-\sum_{n\ge0}\frac{(w-a)^n}{(z-a)^{n+1}},
\]
uniformly, and each term has a primitive, so the integral is $0$. Therefore
\[
\begin{aligned}f(z)&=\frac1{2\pi i}\int_{C(a,\rho_2)}\frac{f(w)\dd w}{w-z}\eqbr\quad-\frac1{2\pi i}\int_{C(a,\rho_1)}\frac{f(w)\dd w}{w-z}.\end{aligned}
\]

\textit{The two series.} Expand the two kernels as geometric series and integrate term by term. The outer integral becomes $\sum_{n\ge0}c_n(z-a)^n$. The inner integral becomes $\sum_{m\ge1}c_{-m}(z-a)^{-m}$. In the variable $u=\frac1{z-a}$, the second series is a power series in $u$, convergent for $|u|<\frac1{\rho_1}$ for every $\rho_1>r$. Lemma~\ref{lem:abel} and Theorem~\ref{thm:powerseries} give its convergence and its holomorphy for $|z-a|>r$. The first series is a power series of radius $R$, by the same two results.

\textit{Uniqueness.} If $f=\sum d_n(z-a)^n$ converges uniformly on $C(a,\rho)$, multiply by $(z-a)^{-k-1}$ and integrate term by term. Proposition~\ref{prop:pathint}(c) returns $2\pi i\,d_k$, so $d_k=c_k$.
\end{proof}

At an isolated singularity one has $r=0$. The \emph{principal part} is $P(z)=\sum_{n<0}c_n(z-a)^n$. It is holomorphic on $\C\setminus\{a\}$, and $f-P$ extends holomorphically through $a$. The singularity is \emph{removable} when $P=0$. It is a \emph{pole of order $m$} when $P$ stops at the term $c_{-m}\ne0$. It is \emph{essential} when $P$ has infinitely many terms. The model essential singularity is $e^{1/z}=\sum_{n\ge0}\frac{z^{-n}}{n!}$ at $0$.

\begin{definition}
The \emph{residue} of $f$ at the isolated singularity $a$ is $\operatorname{Res}_af=c_{-1}=\frac1{2\pi i}\int_{C(a,\rho)}f$ for small $\rho$.
\end{definition}

The residue is the one Laurent coefficient that a small circle can see. Every other coefficient integrates to zero against $1$, by Proposition~\ref{prop:pathint}(c).

\begin{proposition}[computing residues]\label{prop:rescalc}
\begin{enumerate}[label=(\alph*)]
\item If $a$ is a pole of order at most $m$, then $\operatorname{Res}_af=\frac1{(m-1)!}\lim_{z\to a}\frac{\mathrm d^{m-1}}{\mathrm dz^{m-1}}\bigl((z-a)^mf(z)\bigr)$; for a simple pole,
$\operatorname{Res}_af=\lim_{z\to a}(z-a)f(z)$.
\item If $g$, $h$ are holomorphic near $a$ with $h(a)=0\ne h'(a)$, then $\operatorname{Res}_a\frac gh=\frac{g(a)}{h'(a)}$.
\end{enumerate}
\end{proposition}
\begin{proof}
(a) The product $(z-a)^mf(z)=\sum_{n\ge-m}c_n(z-a)^{n+m}$ is a power series. The coefficient of $(z-a)^{m-1}$ in that series is $c_{-1}$. Theorem~\ref{thm:powerseries} reads it off by differentiating $m-1$ times and setting $z=a$. For $m=1$ the formula is the limit $(z-a)f(z)\to c_{-1}$.

(b) The zero of $h$ is simple, so the pole of $g/h$ is at most simple. Then
\[
(z-a)\frac{g(z)}{h(z)}=g(z)\frac{z-a}{h(z)-h(a)}\to\frac{g(a)}{h'(a)}.
\]
\end{proof}

Here is the residue we use for sums. At each integer $n$,
\[
\operatorname{Res}_n(\pi\cot\pi z)=\frac{\pi\cos(\pi n)}{\pi\cos(\pi n)}=1,
\]
by part (b), since the derivative of $\sin\pi z$ at $n$ is $\pi\cos(\pi n)\ne0$. So $\pi\cot\pi z$ has residue $1$ at every integer. That is why it turns a sum over the integers into a contour integral, in the partial-fraction theorem below.

\section{The residue theorem}

For a closed path $\gamma$ and a point $a\notin\gamma^*$, the \emph{winding number} is
\[
n(\gamma,a)=\frac1{2\pi i}\int_\gamma\frac{\mathrm dz}{z-a}.
\]

\begin{proposition}\label{prop:winding}
$n(\gamma,a)$ is an integer: the number of times $\gamma$ turns around $a$.
\end{proposition}
\begin{proof}
Let $\gamma:[\alpha,\beta]\to\C$ and $\phi(t)=\int_\alpha^t\frac{\gamma'(s)}{\gamma(s)-a}\dd s$. Wherever $\gamma$ is $C^1$,
\[
\frac{\mathrm d}{\mathrm dt}\bigl((\gamma(t)-a)e^{-\phi(t)}\bigr)=0.
\]
The continuous function $(\gamma(t)-a)e^{-\phi(t)}$ is therefore constant on $[\alpha,\beta]$. The path is closed, so $\gamma(\beta)=\gamma(\alpha)$ and $e^{\phi(\beta)}=1$. Hence $\phi(\beta)\in2\pi i\Z$. Along the path, $e^{\phi(t)}=\frac{\gamma(t)-a}{\gamma(\alpha)-a}$, so the imaginary part of $\phi$ is a continuous choice of the argument of $\gamma(t)-a$. The integer $\frac{\operatorname{Im}\phi(\beta)}{2\pi}$ counts the full turns.
\end{proof}

Our contours bound bounded convex open sets: discs, half-discs, rectangles, squares, and sectors of angle at most $\pi$. For each point $b$ inside such a set $D$, every ray from $b$ meets the boundary exactly once, at $b+\rho(\theta)e^{i\theta}$. On each side or arc the function $\rho$ is given by an elementary piecewise $C^1$ formula. Call a closed path $\gamma$ the \emph{positive boundary} of $D$ when, for every $b\in D$, the path $\gamma$ is a reparametrisation of $\theta\mapsto b+\rho(\theta)e^{i\theta}$ for $\theta\in[\theta_0,\theta_0+2\pi]$. That is the counterclockwise contour. Then
\[
\begin{aligned}2\pi i\,n(\gamma,b)&=\int_{\theta_0}^{\theta_0+2\pi}\frac{\rho'(\theta)e^{i\theta}+i\rho(\theta)e^{i\theta}}{\rho(\theta)e^{i\theta}}\dd\theta\eqbr=\bigl[\log\rho\bigr]_{\theta_0}^{\theta_0+2\pi}+2\pi i=2\pi i,\end{aligned}
\]
so $n(\gamma,b)=1$ for every $b\in D$. The logarithm of $\rho$ returns to itself after one full turn. The $2\pi i$ is the turn itself.

\begin{theorem}[residue theorem]\label{thm:residue}
Let $D$ be a bounded convex open set with positive boundary $\gamma$, and $U$ an open set containing $\overline D$. Let $f$ be holomorphic on
$U\setminus\{a_1,\dots,a_k\}$, where $a_1,\dots,a_k\in D$. Then
\[
\int_\gamma f(z)\dd z=2\pi i\sum_{j=1}^k\operatorname{Res}_{a_j}f.
\]
\end{theorem}
\begin{proof}
\textit{Remove the principal parts.} Let $P_j$ be the principal part of $f$ at $a_j$. Theorem~\ref{thm:laurent} says $P_j$ is holomorphic on $\C\setminus\{a_j\}$. Near $a_j$, the difference $f-\sum_i P_i$ equals $f-P_j$ minus the other principal parts. The function $f-P_j$ extends holomorphically through $a_j$, and each $P_i$ with $i\ne j$ is holomorphic at $a_j$. So $g=f-\sum_i P_i$ extends to a holomorphic function on $U$.

\textit{The integral of $g$.} For small $\varepsilon>0$, let $D_\varepsilon$ be the set of points at distance less than $\varepsilon$ from $\overline D$. It is open. It is convex: if $|z-d|<\varepsilon$ and $|z'-d'|<\varepsilon$ with $d,d'\in\overline D$, the same inequality holds for convex combinations, because $\overline D$ is convex. For small $\varepsilon$ the set $D_\varepsilon$ lies in $U$, since the compact set $\overline D$ is at positive distance from $\C\setminus U$. Cauchy's theorem on the star-shaped set $D_\varepsilon$ gives $\int_\gamma g=0$.

\textit{Each principal part contributes its residue.} The path $\gamma^*$ stays at some distance $\delta>0$ from $a_j$. On and outside that distance, Theorem~\ref{thm:laurent} gives uniform convergence of $P_j(z)=\sum_{m\ge1}c_{-m}(z-a_j)^{-m}$. Integrate term by term. For $m\ge2$ the term $(z-a_j)^{-m}$ has a primitive, so its integral on the closed path is $0$. The term $m=1$ contributes
\[
c_{-1}\cdot 2\pi i\,n(\gamma,a_j)=2\pi i\operatorname{Res}_{a_j}f,
\]
because $n(\gamma,a_j)=1$.
\end{proof}

\begin{figure}[ht]
\centering
\fig{ch09-contours}
\caption{The three contours of this chapter, with the singularities they enclose. Left: a half-disc of radius $R$, for integrals of
rational functions over $\R$. Middle: a rectangle of height $2\pi$, for integrands periodic under $t\mapsto t+2\pi i$ up to a factor. Right:
the square $C_N$ with vertices $(N+\frac12)(\pm1\pm i)$, which avoids the integers, the poles of $\pi\cot\pi z$.}
\label{fig:contours}
\end{figure}

\section{Real integrals by residues}

The pattern is the same each time. Extend the integrand to the complex plane. Close the segment $[-R,R]$ by a path whose integral is small, or is a known multiple of the integral you want. Let $R\to\infty$.

\begin{proposition}\label{prop:rationalint}
Let $P$, $Q$ be polynomials with $\deg Q\ge\deg P+2$ and no real zero of $Q$. Then
\[
\int_{-\infty}^\infty\frac{P(x)}{Q(x)}\dd x=2\pi i\sum_{\operatorname{Im}a>0}\operatorname{Res}_a\frac PQ,
\]
the sum running over the zeros of $Q$ in the upper half-plane.
\end{proposition}
\begin{proof}
The degree hypothesis gives constants $C$ and $R_0$ with $\bigl|\frac{P(z)}{Q(z)}\bigr|\le\frac C{|z|^2}$ for $|z|\ge R_0$. The real integral converges absolutely. Take $R>R_0$ larger than the modulus of every zero of $Q$. Apply Theorem~\ref{thm:residue} to the half-disc of radius $R$ (Figure~\ref{fig:contours}, left). The integral over $[-R,R]$ plus the integral over the semicircular arc $\Gamma_R$ equals $2\pi i$ times the sum of the residues in the upper half-plane.

The quantity we send to zero is $\int_{\Gamma_R}\frac{P}{Q}$. The arc has length $\pi R$. On it, $\bigl|\frac{P}{Q}\bigr|\le\frac C{R^2}$. Proposition~\ref{prop:pathint}(a) therefore bounds the arc integral by
\[
\pi R\cdot\frac C{R^2}=\frac{\pi C}R.
\]
For this fixed pair of polynomials, $C$ does not depend on $R$. As $R\to\infty$ the bound $\pi C/R$ tends to $0$. The integral over $[-R,R]$ tends to the integral over $\R$, and the arc contributes nothing in the limit.
\end{proof}

\begin{example}
$\int_{-\infty}^\infty\frac{\mathrm dx}{1+x^4}=\frac\pi{\sqrt2}$. The poles in the upper half-plane are $a=e^{i\pi/4}$ and $e^{3i\pi/4}$ (Figure~\ref{fig:rw09-poles}). By Proposition~\ref{prop:rescalc}(b) the
residue at a root $a$ of $z^4=-1$ is $\frac1{4a^3}=\frac a{4a^4}=-\frac a4$, and $e^{i\pi/4}+e^{3i\pi/4}=i\sqrt2$. So the integral is $2\pi i\cdot\bigl(-\frac{i\sqrt2}4\bigr)=\frac{\pi\sqrt2}2=\frac\pi{\sqrt2}$.
\end{example}

The last step is the identification $\frac{\sqrt2}2=\frac1{\sqrt2}$, so $\frac{\pi\sqrt2}2=\frac\pi{\sqrt2}$. Numerically, $\sqrt2\approx1.41421356$ and $\pi\approx3.14159265$, and both $\pi\sqrt2/2$ and $\pi/\sqrt2$ equal $2.22144147$ to eight decimals. The two ways of writing the residue sum are the same number.

\begin{figure}[ht]
\centering
\fig{ch09-rewrite-poles}
\caption{The upper half-plane contour for $\frac1{1+z^4}$. The arc $\Gamma_R$ and the segment $[-R,R]$ enclose the poles at $e^{i\pi/4}$ and $e^{i3\pi/4}$.}
\label{fig:rw09-poles}
\end{figure}

\begin{example}[the transform of $\frac1{1+x^2}$, promised in Chapter~\ref{ch:transform}]
For $\xi\le0$ and $\operatorname{Im}z\ge0$, $|e^{-2\pi iz\xi}|=e^{2\pi\xi\operatorname{Im}z}\le1$. The proof of Proposition~\ref{prop:rationalint} applies to $f(z)=\frac{e^{-2\pi iz\xi}}{1+z^2}$. The quantity sent to zero is the integral of $f$ over the arc $\Gamma_R$. On that arc $|1+z^2|\ge R^2-1$ and the exponential factor has modulus at most $1$, so the arc integral is at most $\frac{\pi R}{R^2-1}$. As $R\to\infty$ this bound tends to $0$, because the denominator grows like $R^2$ and the numerator grows like $R$. The only pole in the upper half-plane is $i$, with residue $\frac{e^{2\pi\xi}}{2i}$. The integral equals $2\pi i$ times that residue, which is $\pi e^{2\pi\xi}=\pi e^{-2\pi|\xi|}$. For $\xi\ge0$ the integrand is even in $\xi$ after the substitution $x\mapsto -x$. This is Exercise~\ref{ex:lorentz}, proved here without the inversion theorem. At $\xi=\frac1{2\pi}$ the real part is the integral from the first page: $\int_{-\infty}^\infty\frac{\cos x}{1+x^2}\dd x=\pi e^{-1}$.
\end{example}

\begin{example}[the reflection formula again]\label{ex:rectangle}
For $0<a<1$, $\displaystyle\int_{-\infty}^\infty\frac{e^{at}}{1+e^t}\dd t=\frac\pi{\sin\pi a}$. The integrand $f$ satisfies $f(t+2\pi i)=e^{2\pi ia}f(t)$. That suggests the rectangle with vertices $\pm R$ and $\pm R+2\pi i$ (Figure~\ref{fig:contours}, middle). Inside, $1+e^t$ vanishes only at $i\pi$, and the residue is $\frac{e^{i\pi a}}{e^{i\pi}}=-e^{i\pi a}$ by Proposition~\ref{prop:rescalc}(b). On the right side $|f|\le\frac{e^{aR}}{e^R-1}$. On the left side $|f|\le\frac{e^{-aR}}{1-e^{-R}}$. Both sides have length $2\pi$, and their integrals tend to $0$ because $0<a<1$. The top side, traversed from right to left, contributes $-e^{2\pi ia}$ times the integral $I_R$ over $[-R,R]$. So
\[
\begin{aligned}(1-e^{2\pi ia})I&=-2\pi ie^{i\pi a},\eqbr I=\frac{2\pi i}{e^{i\pi a}-e^{-i\pi a}}=\frac\pi{\sin\pi a}.\end{aligned}
\]
The substitution $x=e^t$ turns $I$ into $\int_0^\infty\frac{x^{a-1}}{1+x}\dd x$. That integral is $B(a,1-a)=\Gamma(a)\Gamma(1-a)$: substitute $x=\frac u{1-u}$ in the Beta integral of Theorem~\ref{thm:betagamma}. This is a second proof of Theorem~\ref{thm:reflection}, and it does not use the sine product.
\end{example}

The two vertical integrals are the ones that have to vanish, and each bound uses a different half of the condition $0<a<1$. On the right side the quantity being bounded is $\int |f|\,|dz|$. The length is $2\pi$ and $|f|\le\frac{e^{aR}}{e^R-1}$. Since $a<1$, the exponent $(a-1)R$ tends to $-\infty$, so $e^{aR}/(e^R-1)$ tends to $0$ and the product with $2\pi$ tends to $0$. On the left side $|f|\le\frac{e^{-aR}}{1-e^{-R}}$. Since $a>0$, the factor $e^{-aR}$ tends to $0$ while the denominator tends to $1$. The length is again $2\pi$, so that integral tends to $0$ as well.

\begin{example}[the Gaussian, a third time]
For real $\xi$, the function $f(z)=e^{-\pi z^2}$ is entire, so its integral over the rectangle with vertices $\pm R$, $\pm R+i\xi$ is $0$. On the vertical sides, $|e^{-\pi(\pm R+iy)^2}|=e^{-\pi R^2+\pi y^2}\le e^{-\pi R^2+\pi\xi^2}$, so those integrals tend to $0$. Hence $\int_{\R}e^{-\pi(x+i\xi)^2}\dd x=\int_{\R}e^{-\pi x^2}\dd x=1$. Expand $e^{-\pi(x+i\xi)^2}=e^{-\pi x^2}e^{-2\pi ix\xi}e^{\pi\xi^2}$. The result is $\hat g(\xi)=e^{-\pi\xi^2}$. That is Theorem~\ref{thm:gaussianft}, proved by shifting the line of integration.
\end{example}

Name the bound. Each vertical side has length $|\xi|$, or $2|\xi|$ for the two of them together if you prefer one estimate. On either side the integrand has modulus at most $e^{-\pi R^2+\pi\xi^2}$. The product of length and maximum is at most $2|\xi|e^{\pi\xi^2}e^{-\pi R^2}$. For fixed $\xi$, the factor $e^{-\pi R^2}$ tends to $0$ as $R\to\infty$, and everything else stays fixed. The vertical integrals vanish in the limit, and the two horizontal integrals become the identity in the example.

\begin{beauty}
In these three examples the real integral never meets $i$. The value is nevertheless decided at one complex point: $e^{i\pi/4}$ for $\frac1{1+x^4}$, then $i$ for the Fourier transform, then $i\pi$ for the reflection formula. Holomorphy carries the values from the line out into the plane, and the residue is the number in which those values are stored.
\end{beauty}

\section{Sums by residues}

The function $\pi\cot\pi z$ has residue $1$ at every integer. If $f$ is holomorphic at an integer $n$, then $\pi\cot(\pi z)f(z)$ has residue $f(n)$ there. A large square, whose sides miss the integers, turns a sum over the integers into a short list of residues of $f$.

\begin{lemma}\label{lem:cotbound}
On the square $C_N$ with vertices $(N+\frac12)(\pm1\pm i)$, $N\ge0$, we have $|\cot\pi z|\le\coth\frac\pi2<2$.
\end{lemma}
\begin{proof}
Write $\cot\pi z=i\frac{e^{2\pi iz}+1}{e^{2\pi iz}-1}$. On the horizontal sides, $|\operatorname{Im}z|=N+\frac12$, so $q=|e^{2\pi iz}|=e^{-2\pi\operatorname{Im}z}$ is either at most $e^{-\pi}$ or at least $e^\pi$. In both cases, for $|q'|=q$,
\[
\Bigl|\frac{q'+1}{q'-1}\Bigr|\le\frac{q+1}{|q-1|}\le\frac{1+e^{-\pi}}{1-e^{-\pi}}=\coth\frac\pi2.
\]
On the vertical sides, $z=\pm(N+\frac12)+iy$ and $\cot\pi z=-\tan(i\pi y)=-i\tanh\pi y$, whose modulus is less than $1$. Finally $\coth\frac\pi2\approx1.09$, which is less than $2$.
\end{proof}

The bound is uniform in $N$. The square grows, but $|\cot\pi z|$ on it never exceeds $\coth\frac\pi2$. That is why a factor $\frac1{|z|^2}$ will beat the length of the square.

\begin{theorem}\label{thm:residuesum}
Let $f$ be holomorphic on $\C$ except at a finite set $S$, with $|f(z)|\le\frac C{|z|^2}$ for $|z|\ge R_0$. Then
\[
\sum_{n\in\Z\setminus S}f(n)=-\sum_{a\in S}\operatorname{Res}_a\bigl(\pi\cot(\pi z)f(z)\bigr),
\]
the series converging absolutely.
\end{theorem}
\begin{proof}
Absolute convergence of the series follows from the bound on $f$. Let $N$ be large enough that $N+\frac12>R_0$ and that $S$ lies inside $C_N$. The function $g(z)=\pi\cot(\pi z)f(z)$ is holomorphic near the closed square except at the integers $n$ with $|n|\le N$ and at the points of $S$. At an integer $n\notin S$ the pole is simple and the residue is $f(n)$, by Proposition~\ref{prop:rescalc}(a) and the residue $1$ of $\pi\cot\pi z$. Theorem~\ref{thm:residue} gives
\[
\frac1{2\pi i}\int_{C_N}g(z)\dd z=\sum_{|n|\le N,\ n\notin S}f(n)+\sum_{a\in S}\operatorname{Res}_ag.
\]
The quantity sent to zero is the left-hand side. The square has perimeter $8\bigl(N+\frac12\bigr)$. Lemma~\ref{lem:cotbound} gives $|\cot\pi z|\le2$, so $|\pi\cot\pi z|\le2\pi$, and $|f(z)|\le C\big/\bigl(N+\frac12\bigr)^2$ on the square. Proposition~\ref{prop:pathint}(a) bounds the modulus of the normalised integral by
\[
\frac1{2\pi}\cdot 8\Bigl(N+\frac12\Bigr)\cdot 2\pi\cdot\frac C{\bigl(N+\tfrac12\bigr)^2}=\frac{8C}{N+\tfrac12}.
\]
As $N\to\infty$ this bound tends to $0$: the numerator is fixed and the denominator grows like $N$. The finite sum of residues at points of $S$ stays put. What remains is the sum over the integers.
\end{proof}

The Laurent series of $\pi\cot\pi z$ at $0$ starts as follows. From $\sin w=w\bigl(1-\frac{w^2}6+O(w^4)\bigr)$ and $\cos w=1-\frac{w^2}2+O(w^4)$,
\[
w\cot w=\Bigl(1-\frac{w^2}2+\cdots\Bigr)\Bigl(1+\frac{w^2}6+\cdots\Bigr)=1-\frac{w^2}3+O(w^4).
\]
The product is an even series. Substitute $w=\pi z$:
\[
\pi\cot\pi z=\frac1z-\frac{\pi^2}3z+O(z^3).
\]

\begin{corollary}[Basel, the sixth proof]\label{cor:baselres}
$\displaystyle\sum_{n\ge1}\frac1{n^2}=\frac{\pi^2}6$.
\end{corollary}
\begin{proof}
Take $f(z)=\frac1{z^2}$ and $S=\{0\}$. The Laurent series $\frac{\pi\cot\pi z}{z^2}=\frac1{z^3}-\frac{\pi^2}3\cdot\frac1z+O(z)$ has residue $-\frac{\pi^2}3$ at $0$. Theorem~\ref{thm:residuesum} gives $\sum_{n\ne0}\frac1{n^2}=\frac{\pi^2}3$. Half of that sum is the sum over positive $n$.
\end{proof}

\begin{corollary}[Euler's partial fractions, on $\C$]\label{cor:cotc}
For every $z\in\C\setminus\Z$, $\displaystyle\pi\cot\pi z=\frac1z+\sum_{n\ge1}\frac{2z}{z^2-n^2}$.
\end{corollary}
\begin{proof}
Fix $z\notin\Z$. Apply Theorem~\ref{thm:residuesum} to $f(w)=\frac1{w(w-z)}$, in the variable $w$, with $S=\{0,z\}$. At $w=z$ the residue of $\frac{\pi\cot\pi w}{w(w-z)}$ is $\frac{\pi\cot\pi z}z$. At $w=0$, multiply $\pi\cot\pi w=\frac1w+O(w)$ by $\frac1w\cdot\frac1{w-z}=\frac1w\bigl(-\frac1z-\frac w{z^2}+O(w^2)\bigr)$. The coefficient of $\frac1w$ is $-\frac1{z^2}$. So $\sum_{n\ne0}\frac1{n(n-z)}=\frac1{z^2}-\frac{\pi\cot\pi z}z$. Multiply by $-z$ and use $\frac{-z}{n(n-z)}=\frac1{z-n}+\frac1n$. Pair $n$ with $-n$. The terms $\frac1n$ cancel, and $\frac1{z-n}+\frac1{z+n}=\frac{2z}{z^2-n^2}$.
\end{proof}

Chapter~\ref{ch:series} proved the same expansion for real $x$, by Herglotz's trick. Both sides of the identity are now holomorphic on $\C\setminus\Z$, and the residues produce the identity on that set directly. The same method with $f(z)=z^{-2k}$ gives $\zeta(2k)=-\frac12\operatorname{Res}_0\frac{\pi\cot\pi z}{z^{2k}}$, which is half the coefficient of $z^{2k-1}$ in $-\pi\cot\pi z$. The Bernoulli expansion
\[
\pi z\cot\pi z=\sum_{k\ge0}\frac{(-1)^kB_{2k}(2\pi z)^{2k}}{(2k)!}
\]
comes from $\frac t{e^t-1}+\frac t2=\frac t2\coth\frac t2$ at $t=2\pi iz$. Substituting it recovers Euler's formula of Chapter~\ref{ch:series}.

\section{The Gamma function in the complex plane}\label{sec:gammac}

\begin{theorem}\label{thm:gammac}
\begin{enumerate}[label=(\alph*)]
\item $\Gamma(z)=\int_0^\infty t^{z-1}e^{-t}\dd t$ is holomorphic on $\operatorname{Re}z>0$, with $|\Gamma(z)|\le\Gamma(\operatorname{Re}z)$, and $\Gamma(z+1)=z\Gamma(z)$.
\item $\Gamma$ extends to a holomorphic function on the complement of $\{0,-1,-2,\dots\}$, with simple poles at $-n$ of residue $\frac{(-1)^n}{n!}$, satisfying
$\Gamma(z+1)=z\Gamma(z)$ everywhere.
\item $\Gamma(z)\Gamma(1-z)=\frac\pi{\sin\pi z}$ for $z\in\C\setminus\Z$. Hence $\Gamma$ has no zeros, and $\frac1\Gamma$ extends to an entire function vanishing exactly at
$0,-1,-2,\dots$
\item For real $y$, $|\Gamma(\frac12+iy)|^2=\frac\pi{\cosh\pi y}$.
\end{enumerate}
\end{theorem}
\begin{proof}
(a) Here $t^{z-1}=e^{(z-1)\log t}$, so $|t^{z-1}e^{-t}|=t^{x-1}e^{-t}$ with $x=\operatorname{Re}z$. That is the bound $|\Gamma(z)|\le\Gamma(x)$. On a closed disc in the strip $x_1\le\operatorname{Re}z\le x_2$ with $x_1>0$, the integrand is dominated by $t^{x_1-1}e^{-t}+t^{x_2-1}e^{-t}$, which is integrable on $\left]0,\infty\right[$. Theorem~\ref{thm:holoparam}(b) gives holomorphy. Both sides of $\Gamma(z+1)=z\Gamma(z)$ are holomorphic on the connected half-plane $\operatorname{Re}z>0$, and they agree on $\left]0,\infty\right[$ by Proposition~\ref{prop:gammafe}. Theorem~\ref{thm:identity} spreads the equality to the half-plane.

(b) For $n\ge0$, the quotient $\Gamma_n(z)=\frac{\Gamma(z+n+1)}{z(z+1)\cdots(z+n)}$ is holomorphic on $\operatorname{Re}z>-n-1$ except at $0,-1,\dots,-n$, and it equals $\Gamma$ on $\operatorname{Re}z>0$ by the functional equation. Where both $\Gamma_n$ and $\Gamma_{n+1}$ are defined they agree, by that same identity. Together they define the extension. The functional equation holds on the connected set $\C\setminus\{0,-1,\dots\}$ by the identity theorem. Near $-n$,
\[
(z+n)\Gamma(z)=\frac{\Gamma(z+n+1)}{z(z+1)\cdots(z+n-1)}
\]
tends to $\frac{\Gamma(1)}{(-n)(-n+1)\cdots(-1)}=\frac{(-1)^n}{n!}$. The pole is simple, and that limit is the residue.

(c) Both sides of the reflection formula are holomorphic on $\C\setminus\Z$, which is connected, and they agree on $\left]0,1\right[$ by Theorem~\ref{thm:reflection}. If $\Gamma(z_0)=0$, then $z_0\notin\Z$, because $\Gamma(n)=(n-1)!$ at the positive integers. The left side would vanish at $z_0$ and the right side does not. So $\Gamma$ has no zeros. The reciprocal $\frac1\Gamma$ is holomorphic off the poles of $\Gamma$, and at each pole it tends to $0$. Theorem~\ref{thm:removable} makes it entire.

(d) Conjugating the integral gives $\overline{\Gamma(z)}=\Gamma(\overline z)$. Therefore
\[
\bigl|\Gamma(\tfrac12+iy)\bigr|^2=\Gamma(\tfrac12+iy)\Gamma(\tfrac12-iy)=\frac\pi{\sin(\frac\pi2+i\pi y)}=\frac\pi{\cos(i\pi y)}=\frac\pi{\cosh\pi y}.
\]
\end{proof}

On the line $\operatorname{Re}z=\frac12$, part (d) gives $|\Gamma(\frac12+iy)|^2=\pi/\cosh(\pi y)$. For large $|y|$, $\cosh(\pi y)$ is about $\frac12 e^{\pi|y|}$, so $|\Gamma(\frac12+iy)|$ is about $\sqrt{2\pi}\,e^{-\pi|y|/2}$. That vertical decay is why $\Gamma$ is a useful factor in Chapter~\ref{ch:zeta}: the product $\pi^{-s/2}\Gamma(\frac s2)\zeta(s)$ can be studied on vertical lines. Every identity of Chapter~\ref{ch:gamma} between functions that extend holomorphically, the duplication formula for instance, extends to the complex plane by the identity theorem, exactly as in (a). Exercise~\ref{ex:wielandt} is the complex stand-in for Bohr--Mollerup: $\Gamma$ is the only solution of $F(z+1)=zF(z)$, $F(1)=1$, that stays bounded on a vertical strip.

\section{Counting zeros}\label{sec:argument}

A function is \emph{meromorphic} on $U$ when it is holomorphic except at isolated poles.

\begin{theorem}[argument principle]\label{thm:argument}
Let $D$ be a bounded convex open set with positive boundary $\gamma$, and $f$ meromorphic on an open set $U\supseteq\overline D$, with finitely many poles,
not identically zero on any disc, and with no zeros or poles on $\gamma^*$. Then
\[
\frac1{2\pi i}\int_\gamma\frac{f'(z)}{f(z)}\dd z=Z-P,
\]
where $Z$ and $P$ are the numbers of zeros and poles of $f$ in $D$, counted with multiplicity. The left side is also $n(f\circ\gamma,0)$.
\end{theorem}
\begin{proof}
Shrink $U$ to the convex set $D_\varepsilon$ from the proof of Theorem~\ref{thm:residue}. That set is connected. The zeros of $f$ in the compact set $\overline D$ are finite in number, because they are isolated off the poles (Theorem~\ref{thm:identity}). At a zero of order $m$, one has $f=(z-a)^mg$ with $g(a)\ne0$, so
\[
\frac{f'}f=\frac m{z-a}+\frac{g'}g
\]
has residue $m$. At a pole of order $m$, one has $f=(z-a)^{-m}g$ with $g$ holomorphic and $g(a)\ne0$, and the residue is $-m$. Elsewhere $\frac{f'}f$ is holomorphic. Theorem~\ref{thm:residue} adds the residues. The chain rule gives the winding-number form:
\[
\int_{f\circ\gamma}\frac{\mathrm dw}w=\int_\alpha^\beta\frac{f'(\gamma(t))\gamma'(t)}{f(\gamma(t))}\dd t.
\]
\end{proof}

The number of zeros inside a contour is the number of times the image of the contour winds around $0$ (Figure~\ref{fig:winding}). For a polynomial of degree $n$ on a huge circle, $f(z)$ is close to $z^n$, and the image winds $n$ times. That is the fundamental theorem of algebra again, with the count included. Rouch\'e's theorem makes the word ``close'' into a hypothesis you can check.

\begin{theorem}[Rouch\'e]\label{thm:rouche}
Let $D$, $\gamma$ be as above, $f$, $g$ holomorphic on an open set containing $\overline D$, and $|g|<|f|$ on $\gamma^*$. Then $f$ and $f+g$ have the same
number of zeros in $D$.
\end{theorem}
\begin{proof}
For $s\in[0,1]$, set $f_s=f+sg$. On $\gamma^*$ one has $|f_s|\ge|f|-|g|>0$, so $f_s$ has no zero on the contour and is not identically zero. The argument principle applies. The number of zeros in $D$ is
\[
N(s)=\frac1{2\pi i}\int_\gamma\frac{f'+sg'}{f+sg}.
\]
The integrand is continuous in $(s,t)$ on the compact set $[0,1]\times[\alpha,\beta]$, piece by piece. Its denominator is bounded below by $\min_{\gamma^*}(|f|-|g|)>0$. Theorem~\ref{thm:paramcompact} says $N$ is continuous. A continuous integer-valued function on $[0,1]$ is constant, so $N(0)=N(1)$.
\end{proof}

\begin{example}
How many roots of $z^5+3z+1$ lie in the unit disc? On $|z|=1$, one has $|z^5+1|\le2<3=|3z|$. Rouch\'e's theorem with $f=3z$ and $g=z^5+1$ says there is exactly one, the same count as for $3z$. On $|z|=\frac32$, one has $|3z+1|\le\frac{11}2<\frac{243}{32}=|z^5|$, so all five roots have modulus less than $\frac32$.
\end{example}

\begin{figure}[ht]
\centering
\fig{ch09-winding}
\caption{Left: the five zeros of $z^5+3z+1$ (rust) and the circles $|z|=1$ (charcoal) and $|z|=\frac32$ (coral). Middle: the image of the unit circle
under $z\mapsto z^5+3z+1$, rescaled, winds once around $0$ (the black dot), despite its four loops. Right: the image of the larger circle winds
five times. Only one zero is inside the unit circle, and all five are inside the larger one.}
\label{fig:winding}
\end{figure}

\begin{think}
The argument principle turns a question about where the zeros are into an integral. Is this a practical way to count roots? What would you do if a numerical quadrature returned $2.9999$?
\end{think}

The integral equals an integer, $Z-P$. A quadrature that returns $2.9999$ is reporting the integer $3$, with an error in the quadrature. You round to the nearest integer once the error is smaller than $\frac12$. In practice you track the continuous change in $\arg f(\gamma(t))$ as $t$ runs along the contour, which is the same integer. Often you do not integrate at all. Rouch\'e's comparison, as in the example of $z^5+3z+1$, counts the roots by an inequality on the circle.

\section{What we proved, and what comes next}

One fact generated the chapter: the integral of a holomorphic function around a triangle is zero (Theorem~\ref{thm:goursat}). From it, in order, we built a primitive on a star-shaped set, with $F'=f$ (Theorem~\ref{thm:cauchy}), Cauchy's formula (Theorem~\ref{thm:cauchyformula}), and the Taylor series on the largest disc that fits (Theorem~\ref{thm:analytic}). That series is why $\frac1{1+x^2}$ stops converging at distance $1$. Liouville, the identity theorem, and the maximum principle came next. Around an isolated singularity the Laurent series (Theorem~\ref{thm:laurent}) has one coefficient, the residue, that a contour integral remembers (Theorem~\ref{thm:residue}).

The residue theorem then ran four times. Real integrals: the integral on the first page is $2\pi i$ times the residue at $i$, and $\int_{-\infty}^\infty\frac{\mathrm dx}{1+x^4}=\frac\pi{\sqrt2}$. Sums: $\pi\cot\pi z$ turns a sum over the integers into a few residues (Theorem~\ref{thm:residuesum}), which gave Basel a sixth proof and Euler's partial fractions on all of $\C$. Extension: $\Gamma$ lives on the plane, with poles at $0,-1,-2,\dots$ and with no zeros (Theorem~\ref{thm:gammac}). Counting: the zeros inside a contour are an integral (Theorem~\ref{thm:argument}), and they do not move when the function is perturbed by something smaller on the contour (Theorem~\ref{thm:rouche}).

Chapter~\ref{ch:zeta} uses all four at once, on the single function $\zeta$. The question there is what the symmetry $s\leftrightarrow 1-s$ says about the values $\zeta(-1)$, $\zeta(0)$, and the zeros. Chapter~\ref{ch:asymptotics} reads the growth of coefficients from the nearest singularity, which Exercise~\ref{pb:zigzagpole} already does for the secant series. Chapter~\ref{ch:hr} integrates over a circle that moves toward a wall of singularities.

Homology, the Phragm\'en--Lindel\"of principle, Hadamard's factorisation, and Mellin inversion remain in Volume~8, Chapter~14. The Jordan curve theorem remains in Volume~7, Chapter~6.

\section{Exercises}

\begin{exercise}[$\star\star$]\label{ex:xn}
For $n\ge2$, prove $\displaystyle\int_0^\infty\frac{\mathrm dx}{1+x^n}=\frac{\pi/n}{\sin(\pi/n)}$, using the sector $0\le\arg z\le\frac{2\pi}n$, $|z|\le R$. The integral over the arc tends to $0$ as $R\to\infty$, for this fixed $n\ge2$. Check the result against the Beta function.
\end{exercise}

\begin{exercise}[$\star\star$]\label{ex:fresnelsector}
Integrate $e^{iz^2}$ over the sector $0\le\arg z\le\frac\pi4$, $|z|\le R$. Prove $\displaystyle\int_0^\infty\cos(x^2)\dd x=\int_0^\infty\sin(x^2)\dd x=\frac12\sqrt{\frac\pi2}$. Compare the result with Lemma~\ref{lem:fresnel}.
\end{exercise}

\begin{exercise}[$\star$]\label{ex:cothres}
Use Theorem~\ref{thm:residuesum} to prove $\displaystyle\sum_{n\ge1}\frac1{n^2+a^2}=\frac\pi{2a}\coth\pi a-\frac1{2a^2}$ for $a>0$. Compare the sum with Exercise~\ref{ex:coth}. What is $\sum_{n\in\Z}\frac{(-1)^n}{n^2+a^2}$?
\end{exercise}

\begin{exercise}[$\star$]\label{ex:rouche}
Prove that $z^4-6z+3$ has exactly one zero in the unit disc and three zeros in the annulus $1<|z|<2$.
\end{exercise}

\begin{exercise}[$\star\star$]\label{ex:polygrowth}
Let $f$ be entire and satisfy $|f(z)|\le C(1+|z|)^k$ for every $z$. Prove that $f$ is a polynomial of degree at most $k$. Deduce that an entire function $f$ with $|f(z)|\le C|e^z|$ for every $z$ has the form $ce^z$.
\end{exercise}

\begin{exercise}[$\star\star\star$]\label{ex:wielandt}
(Wielandt's theorem, 1939.) Let $F$ be holomorphic on $\operatorname{Re}z>0$, with $F(z+1)=zF(z)$, $F(1)=1$, and with $F$ bounded on the strip $1\le\operatorname{Re}z\le2$. Prove $F=\Gamma$.
\begin{enumerate}[label=(\alph*)]
\item Prove that $\Gamma$ is bounded on the strip, and set $G=F-\Gamma$. Prove that $G$ extends to an entire function satisfying $G(z+1)=zG(z)$.
\item Prove that $H(z)=G(z)G(1-z)$ is bounded on $0\le\operatorname{Re}z\le1$ and satisfies $H(z+1)=-H(z)$. Finish the argument.
\end{enumerate}
\end{exercise}

\begin{exercise}[$\star\star$]\label{pb:rankdist}
\emph{The rank of a random matrix.}
\begin{enumerate}[label=(\alph*)]
\item Prove that the number of matrices of rank $r$ in $M_n(\mathbb F_q)$ is
\[
\prod_{i=0}^{r-1}\frac{(q^n-q^i)^2}{q^r-q^i}.
\]
\item Let $A$ be uniform in $M_n(\mathbb F_q)$. Prove that for each $k\ge0$, as $n\to\infty$,
\[
P(\operatorname{corank}A=k)\longrightarrow q^{-k^2}\,\frac{\prod_{j>k}(1-q^{-j})}{\prod_{j=1}^k(1-q^{-j})}.
\]
\item Deduce Durfee's identity. For $|x|<1$, write $(x)_k=(1-x)(1-x^2)\cdots(1-x^k)$ and prove
\[
\sum_{k\ge0}\frac{x^{k^2}}{(x)_k^{\,2}}=\prod_{j\ge1}\frac1{1-x^j},
\]
first at $x=\frac1q$, then for every $x$.
\end{enumerate}
\end{exercise}

\begin{exercise}[Journey]\label{pb:dedekind}
\emph{Dedekind sums.} Let $((t))=t-\lfloor t\rfloor-\frac12$ for $t\notin\Z$, and let $((t))=0$ for $t\in\Z$. For coprime integers $h$ and $k\ge1$, the \emph{Dedekind sum} is $s(h,k)=\sum_{r=0}^{k-1}\bigl(\bigl(\frac rk\bigr)\bigr)\bigl(\bigl(\frac{hr}k\bigr)\bigr)$.
\begin{enumerate}[label=(\alph*)]
\item Prove $\bigl(\bigl(\frac rk\bigr)\bigr)=-\frac1{2k}\sum_{j=1}^{k-1}\cot\frac{\pi j}k\,\sin\frac{2\pi jr}k$ for every integer $r$, and deduce
\[
s(h,k)=\frac1{4k}\sum_{j=1}^{k-1}\cot\frac{\pi j}k\cot\frac{\pi hj}k.
\]
\item Let $h,k\ge1$ be coprime. Integrate $\cot(\pi z)\cot(\pi hz)\cot(\pi kz)$ around the rectangle with vertices $\varepsilon\pm iY$, $1+\varepsilon\pm iY$, and let $Y\to\infty$. Prove \emph{Dedekind's reciprocity law}
\[
s(h,k)+s(k,h)=-\frac14+\frac1{12}\Bigl(\frac hk+\frac kh+\frac1{hk}\Bigr).
\]
\item Prove $s(1,k)=\frac{(k-1)(k-2)}{12k}$, prove that $6k\,s(h,k)$ is an integer, and prove $|s(h,k)|\le s(1,k)$.
\item Let $\frac hk=[0;a_1,\dots,a_n]$ with $0<h<k$. Prove $|s(h,k)|\le\frac1{12}(a_1+\dots+a_n)+n$, and prove that $s(h,k)$ can be computed in $O(\log k)$ arithmetic operations on fractions.
\item For $k=101$, $1009$, $10\,007$, $100\,003$, the average of $|s(h,k)|$ over $1\le h<k$ is $0.951$, $2.037$, $3.691$, $5.802$, while $\log^2k$ is $21.3$, $47.8$, $84.8$, $132.5$, and the maximum is $s(1,k)$. Vardi proved in 1993 that $\frac{s(h,k)}{\log k}$ has a Cauchy limiting distribution as $h$ and $k$ vary. Explain why a law with no mean can have averages that grow like $\log^2k$, and not like $\log k$.
\end{enumerate}
\end{exercise}

\begin{exercise}[$\star\star$]\label{pb:zigzagpole}
\emph{Why $2(2/\pi)^{n+1}$?} Let $A_n$ be the number of alternating permutations of $\{1,\dots,n\}$, the ones with $\sigma(1)<\sigma(2)>\sigma(3)<\cdots$. Exercise~\ref{pb:zigzag} of Chapter~\ref{ch:generating} shows that $\sum_nA_n\frac{z^n}{n!}$ is the Taylor series at $0$ of $g(z)=\sec z+\tan z=\frac{1+\sin z}{\cos z}$.
\begin{enumerate}[label=(\alph*)]
\item Prove that $g(z)-\frac4{\pi-2z}$ extends to a holomorphic function on the disc $|z|<\frac{3\pi}2$.
\item Deduce
\[
\frac{A_n}{n!}=2\Bigl(\frac2\pi\Bigr)^{n+1}\bigl(1+O(\delta^n)\bigr)\qq\text{for every }\delta>\tfrac13.
\]
\item For $n=10$, $A_{10}=50521$. Compute the relative error of the approximation, and compare it with $-3^{-11}$.
\end{enumerate}
\end{exercise}

\begin{exercise}[$\star\star\star$]\label{ex:jensen}
\emph{Jensen's formula: zeros cost growth.} Let $f$ be holomorphic on an open set containing $\overline D(0,R)$, with $f(0)\ne0$ and with no zero on the circle $|z|=R$. Let $a_1,\dots,a_n$ be its zeros in $D(0,R)$, repeated according to multiplicity.
\begin{enumerate}[label=(\alph*)]
\item Let $g$ be holomorphic and without zeros on a disc $D(0,R')$. Prove that $g=e^L$ for some holomorphic $L$ on $D(0,R')$, and deduce $\frac1{2\pi}\int_0^{2\pi}\log|g(re^{it})|\dd t=\log|g(0)|$ for $0<r<R'$.
\item For $0<|a|<R$ set $B_a(z)=\frac{R(z-a)}{R^2-\overline az}$. Prove $|B_a(z)|=1$ for $|z|=R$, and compute $|B_a(0)|$.
\item Prove Jensen's formula:
\[
\frac1{2\pi}\int_0^{2\pi}\log\bigl|f(Re^{it})\bigr|\dd t=\log|f(0)|+\sum_{k=1}^n\log\frac R{|a_k|}.
\]
\item Let $f$ be entire, with $f(0)=1$ and $|f(z)|\le Ce^{A|z|}$ for every $z$. Prove that the number $N(r)$ of zeros of $f$ in $\overline D(0,r)$, counted with multiplicity, satisfies $N(r)\le eAr+\log C$. Compare the bound with $f(z)=\frac{\sin\pi z}{\pi z}$.
\end{enumerate}
\end{exercise}

\begin{exercise}[$\star$]\label{ex:rw09a}
On the unit circle, compute $\int_{C(0,1)}z^2\dd z$ and $\int_{C(0,1)}\frac{\mathrm dz}z$ by writing $z=e^{it}$. Check that the second integral equals $2\pi i$, and that $2\pi\approx6.2832$.
\end{exercise}

\begin{exercise}[$\star\star$]\label{ex:rw09b}
Show that $\displaystyle\int_0^{2\pi}\frac{\mathrm d\theta}{2+\cos\theta}=\frac{2\pi}{\sqrt3}$. Use the unit circle and the substitution $z=e^{i\theta}$. Check the value $2\pi/\sqrt3\approx3.6276$ against the $8$-point sum $\frac{2\pi}8\sum_{k=0}^{7}\frac1{2+\cos(2\pi k/8)}$.
\end{exercise}

\begin{exercise}[$\star\star$]\label{ex:rw09c}
Show that $\displaystyle\int_{-\infty}^\infty\frac{\mathrm dx}{(1+x^2)^2}=\frac\pi2$, by a residue of order $2$ at $z=i$. Check the same value from the antiderivative $\frac12\Bigl(\frac x{1+x^2}+\arctan x\Bigr)$.
\end{exercise}

\begin{exercise}[$\star\star\star$]\label{ex:rw09d}
The slit plane is star-shaped from $1$. Theorem~\ref{thm:cauchy} says that $F(z)=\int_{[1,z]}\frac{\mathrm dw}w$ satisfies $F'=1/z$. Parametrise the segment from $1$ to $i$ and show that the integral equals $\operatorname{Log} i=i\pi/2$. As a numerical check, the midpoint rule with four panels on that segment returns about $1.581i$, next to $i\pi/2\approx1.5708i$.
\end{exercise}
\section{Hints}

\begin{hint}{ex:xn}
The only pole in the sector is $e^{i\pi/n}$. On the second ray $z=re^{2\pi i/n}$ you have $z^n=r^n$. For the check, substitute $u=x^n$ and use the reflection formula.
\end{hint}
\begin{hint}{ex:fresnelsector}
On the arc, $|e^{iz^2}|=e^{-R^2\sin2\theta}$, and $\sin2\theta\ge\frac{4\theta}\pi$ on $[0,\frac\pi4]$ because sine is concave there. On the ray $z=re^{i\pi/4}$ you have $iz^2=-r^2$.
\end{hint}
\begin{hint}{ex:cothres}
Compute the residues of $\frac{\pi\cot\pi z}{z^2+a^2}$ at $\pm ia$, and use $\cot(iy)=-i\coth y$. For the alternating sum, $\frac\pi{\sin\pi z}$ has residue $(-1)^n$ at each integer $n$.
\end{hint}
\begin{hint}{ex:rouche}
On $|z|=1$, compare with $-6z$. On $|z|=2$, compare with $z^4$. The strict inequality also shows there is no zero on $|z|=1$.
\end{hint}
\begin{hint}{ex:polygrowth}
Apply Cauchy's estimates on $C(0,r)$, then let $r\to\infty$. For the second claim, look at $f(z)e^{-z}$.
\end{hint}
\begin{hint}{ex:wielandt}
(a) Log-convexity and $\Gamma(1)=\Gamma(2)=1$ bound $\Gamma$ on the strip. Near $-n$, write $G(z)=\frac{G(z+n+1)}{z(z+1)\cdots(z+n)}$ and use $G(1)=0$. (b) For $|\operatorname{Im}z|\ge1$ in the strip, $G(z)=\frac{G(z+1)}z$. Then Liouville and Theorem~\ref{thm:identity} finish the argument.
\end{hint}

\begin{hint}{pb:rankdist}
(a) Choose the image first. (b) Pull $q^{n}$ out of every factor. (c) The probabilities add up to $1$. Pass to the limit with a bound that does not depend on $n$.
\end{hint}

\begin{hint}{pb:dedekind}
(a) The sawtooth on $\Z/k\Z$ has discrete Fourier coefficients $\frac i{2k}\cot\frac{\pi j}k$. Parseval turns the sum into cotangents. (b) Inside the rectangle the poles are a triple pole at $z=1$ and simple poles at $\frac jh$ and $\frac jk$. As $Y\to\infty$, the cotangents tend to $\mp i$ on the horizontal sides. (c) Cauchy--Schwarz gives $|s(h,k)|\le s(1,k)$. (d) Run reciprocity along Euclid's algorithm.
\end{hint}

\begin{hint}{pb:zigzagpole}
(a) The zeros of $\cos$ nearest to $0$ are $\pm\frac\pi2$. At $-\frac\pi2$ the numerator $1+\sin z$ vanishes as well. (b) Cauchy's estimates (Corollary~\ref{cor:radius}) on a circle of radius just below $\frac{3\pi}2$ bound the remainder.
\end{hint}

\begin{hint}{ex:jensen}
(a) The quotient $\frac{g'}g$ has a primitive on the disc, by Theorem~\ref{thm:cauchy}. Differentiate $ge^{-L}$. Then take the real part of the mean value property of $L$. (b) On the circle, $R^2=z\overline z$. (c) Divide $f$ by the product of the factors $B_{a_k}$. (d) Apply Jensen on a circle of radius $R=er$. Each zero in $\overline D(0,r)$ contributes at least $\log\frac Rr=1$.
\end{hint}

\begin{hint}{ex:rw09a}
Write $z=e^{it}$ and $\mathrm dz=ie^{it}\,\mathrm dt$. The integral of $e^{3it}$ over a full period is $0$. The integral of the constant $i$ is $2\pi i$.
\end{hint}
\begin{hint}{ex:rw09b}
Use $\cos\theta=(z+1/z)/2$ and $\mathrm d\theta=\mathrm dz/(iz)$. The pole inside the unit circle is $z=-2+\sqrt3$. The eight-point sum is a check of the number, not the proof.
\end{hint}
\begin{hint}{ex:rw09c}
The pole at $i$ has order $2$. Differentiate $(z+i)^{-2}$ and evaluate at $z=i$. The antiderivative tends to $\pi/4$ at $+\infty$ and to $-\pi/4$ at $-\infty$.
\end{hint}
\begin{hint}{ex:rw09d}
Parametrise the segment by $w=(1-t)+it$, $t\in[0,1]$. The imaginary part of the integrand is $1/(2t^2-2t+1)$, and that integral is $\arctan(2t-1)$ from $0$ to $1$.
\end{hint}

\section{Solutions}

\begin{solution}{ex:xn}
\idea Rotating the ray by $\omega=e^{2\pi i/n}$ leaves the integrand $1/(1+z^n)$ unchanged, because $(\omega z)^n=z^n$. The integral along the rotated ray is $\omega$ times the integral you want. The same comparison appeared on the rectangle of Example~\ref{ex:rectangle}.
\writeup The sector is convex, because its angle $\frac{2\pi}n$ is at most $\pi$. The zeros of $1+z^n$ are $e^{i\pi(2k+1)/n}$, and only $a=e^{i\pi/n}$ has argument in
$\left]0,\frac{2\pi}n\right[$. The residue there is $\frac1{na^{n-1}}=\frac a{na^n}=-\frac an$. The quantity sent to zero is the integral of $\frac1{1+z^n}$ over the arc $|z|=R$. On that arc $\bigl|\frac1{1+z^n}\bigr|\le\frac1{R^n-1}$, and the arc has length $\frac{2\pi R}n$, so the integral is bounded by
\[
\frac{2\pi R/n}{R^n-1}.
\]
Fix $n\ge2$. The numerator grows like $R$ and the denominator grows like $R^n$. For this fixed $n$ the bound tends to $0$ as $R\to\infty$. On the ray $z=r\omega$, traversed inwards, the integral is $-\omega\int_0^R\frac{\mathrm dr}{1+r^n}$. Letting $R\to\infty$, with $\omega=a^2$,
$(1-a^2)I=-\frac{2\pi ia}n$, so
\[
\begin{aligned}I&=\frac{2\pi i}{n(a-a^{-1})}\eqbr=\frac{2\pi i}{n\cdot2i\sin\frac\pi n}=\frac{\pi/n}{\sin(\pi/n)}.\end{aligned}
\]
Check: with $u=x^n$, $I=\frac1n\int_0^\infty\frac{u^{1/n-1}}{1+u}\dd u=\frac1n\Gamma(\frac1n)\Gamma(1-\frac1n)$ by Example~\ref{ex:rectangle}, and reflection gives the same value. For $n=2$, $\frac\pi2$.
\end{solution}

\begin{solution}{ex:fresnelsector}
\idea On the real axis, $e^{ix^2}$ oscillates. On the diagonal $z=re^{i\pi/4}$, the same function is the Gaussian $e^{-r^2}$. Cauchy's theorem moves the integral from one ray to the other once the arc between them is small.
\writeup $e^{iz^2}$ is entire, so its integral over the boundary of the sector is $0$. The sector is convex. On the arc $z=Re^{i\theta}$, one has $iz^2=iR^2(\cos2\theta+i\sin2\theta)$, so
$|e^{iz^2}|=e^{-R^2\sin2\theta}\le e^{-4R^2\theta/\pi}$. The inequality uses that the concave function $\sin2\theta$ lies above its chord on $[0,\frac\pi4]$. The quantity being bounded is the integral of $e^{iz^2}$ over that arc. It is at most
\[
\int_0^{\pi/4}e^{-4R^2\theta/\pi}R\dd\theta\le\frac{\pi R}{4R^2}=\frac\pi{4R}.
\]
For the fixed angle $\frac\pi4$, the bound $\pi/(4R)$ tends to $0$ as $R\to\infty$. On the ray, $z=re^{i\pi/4}$ and $iz^2=-r^2$, contributing $-e^{i\pi/4}\int_0^Re^{-r^2}\dd r$. So
\[
\int_0^Re^{ix^2}\dd x\longrightarrow e^{i\pi/4}\int_0^\infty e^{-r^2}\dd r=\frac{\sqrt\pi}2\cdot\frac{1+i}{\sqrt2},
\]
which proves that the improper integral converges, and taking real and imaginary parts gives $\frac{\sqrt\pi}{2\sqrt2}=\frac12\sqrt{\frac\pi2}$ for both. With $u=x^2$,
$\int_0^\infty e^{ix^2}\dd x=\frac12\int_0^\infty u^{-1/2}e^{iu}\dd u$, which is $\frac12\sqrt\pi e^{i\pi/4}$ by Lemma~\ref{lem:fresnel}: the same value, by a different road.
\end{solution}

\begin{solution}{ex:cothres}
\idea Use the theorem at the two poles $\pm ia$ of $f$. The cotangent of a purely imaginary number is a hyperbolic cotangent.
\writeup The function $f(z)=\frac1{z^2+a^2}$ satisfies the bound, and it has simple poles at $\pm ia$, which are not integers. Proposition~\ref{prop:rescalc}(b) says the residue of $\frac{\pi\cot\pi z}{z^2+a^2}$ at $\pm ia$ is $\frac{\pi\cot(\pm i\pi a)}{\pm2ia}$. Since $\cot(iy)=\frac{\cos iy}{\sin iy}=\frac{\cosh y}{i\sinh y}=-i\coth y$ and $\cot$ is odd, both residues equal $\frac{-i\pi\coth\pi a}{2ia}=-\frac{\pi\coth\pi a}{2a}$. So $\sum_{n\in\Z}\frac1{n^2+a^2}=\frac\pi a\coth\pi a$, the same sum as in Exercise~\ref{ex:coth}, there by Poisson summation. Remove $n=0$ and halve, and the formula follows. For the alternating sum, $\frac\pi{\sin\pi z}$ obeys the same kind of bound as Lemma~\ref{lem:cotbound} on $C_N$: on the horizontal sides $\bigl|\frac{2i}{e^{i\pi z}-e^{-i\pi z}}\bigr|\le\frac2{e^{\pi/2}-e^{-\pi/2}}$, and on the vertical sides $|\sin\pi z|=\cosh\pi y\ge1$. Its residue at $n$ is $(-1)^n$. The same proof gives $\sum_n\frac{(-1)^n}{n^2+a^2}=-2\cdot\frac{\pi}{2ia\sin(i\pi a)}=\frac\pi{a\sinh\pi a}$, using $\sin(iy)=i\sinh y$.
\end{solution}

\begin{solution}{ex:rouche}
\idea On each circle, find the term that is larger than the rest.
\writeup On $|z|=1$, $|z^4+3|\le4<6=|-6z|$, so $z^4-6z+3$ has as many zeros in the unit disc as $-6z$, namely one. The inequality is strict, so there is no zero on the circle. On $|z|=2$, $|-6z+3|\le15<16=|z^4|$, so there are four zeros in $|z|<2$, the same count as for $z^4$. Hence three zeros lie in $1<|z|<2$.
\end{solution}

\begin{solution}{ex:polygrowth}
\idea The proof of Liouville still works for polynomial growth. Cauchy's estimates force every coefficient past the degree to vanish.
\writeup Theorem~\ref{thm:analytic} gives $f(z)=\sum c_nz^n$ on $\C$, and Corollary~\ref{cor:radius}(a) gives $|c_n|\le\frac{C(1+r)^k}{r^n}$ for every $r>0$. For $n>k$ the right side tends to $0$ as $r\to\infty$, so $c_n=0$. For the second claim, $g(z)=f(z)e^{-z}$ is entire and bounded by $C$, so Liouville says it is constant.
\end{solution}

\begin{solution}{ex:wielandt}
\idea On the real axis, Bohr--Mollerup used log-convexity. In the complex plane, boundedness on a strip takes that role, and Liouville finishes the argument. The product $G(z)G(1-z)$ is built so that it is periodic, in the same way that $\Gamma(z)\Gamma(1-z)$ is periodic.
\writeup (a) The function $\Gamma$ is log-convex on $\left]0,\infty\right[$ (Corollary~\ref{cor:gammalogconvex}) with $\Gamma(1)=\Gamma(2)=1$, so $\Gamma\le1$ on $[1,2]$. Theorem~\ref{thm:gammac}(a) gives $|\Gamma(z)|\le\Gamma(\operatorname{Re}z)\le1$ on the strip. So $G$ is bounded on the strip, $G(z+1)=zG(z)$ and $G(1)=0$. As in Theorem~\ref{thm:gammac}(b), the formula $G(z)=\frac{G(z+n+1)}{z(z+1)\cdots(z+n)}$ extends $G$ to $\C\setminus\{0,-1,\dots\}$. Near $-n$, the numerator tends to $G(1)=0$, so $(z+n)G(z)\to0$. The function $G$ is bounded near $-n$: the numerator is holomorphic and vanishes at $-n$, so Theorem~\ref{thm:identity}(a) writes it as $(z+n)$ times a holomorphic function. Theorem~\ref{thm:removable} says $G$ is entire.

(b) On the strip $S_0=\{0\le\operatorname{Re}z\le1\}$, the quotient $G(z)=\frac{G(z+1)}z$ is bounded for $|\operatorname{Im}z|\ge1$, and continuity bounds $G$ on the compact part $|\operatorname{Im}z|\le1$. The map $z\mapsto1-z$ sends $S_0$ to itself, so $H$ is bounded on $S_0$. Next, $G(1-z)=G(-z+1)=-zG(-z)$, so for $z\ne0$,
\[
\begin{aligned}H(z+1)&=G(z+1)G(-z)=zG(z)\cdot\frac{G(1-z)}{-z}\eqbr=-H(z),\end{aligned}
\]
and continuity extends the identity to $z=0$. So $|H|$ is $1$-periodic, hence bounded on $\C$, and Liouville says $H$ is constant: $H=H(1)=G(1)G(0)=0$. If $G$ were not identically zero, its zeros would be isolated, so $G\ne0$ on some disc $D$. Then $G(1-z)=0$ for $z\in D$, and Theorem~\ref{thm:identity} would force $G=0$, a contradiction. Hence $F=\Gamma$.
\end{solution}

\begin{solution}{pb:rankdist}
\idea A matrix of rank $r$ is a choice of an $r$-dimensional image, together with a surjection onto that image. In the limit only the corank matters. Adding the limit probabilities to $1$ is a partition identity.

\writeup (a) A matrix $A$ of rank $r$ is a linear map $\mathbb F_q^n\to\mathbb F_q^n$ whose image $W$ has dimension $r$. Fix $W$ and a basis of $W$. The maps onto $W$ correspond to $r\times n$ matrices of rank $r$, with independent rows: there are $\prod_{i<r}(q^n-q^i)$ of them. The number of choices of $W$ is the Gaussian binomial $\prod_{i<r}\frac{q^n-q^i}{q^r-q^i}$: count ordered independent vectors, then divide by the number of bases of an $r$-dimensional space. Multiply the two counts.

(b) Set $r=n-k$. Factor $q^n$ out of each $q^n-q^i$, and factor $q^r$ out of each $q^r-q^i$:
\[
\frac{\#\{\text{rank }r\}}{q^{n^2}}=q^{2nr-r^2-n^2}\,\frac{\prod_{i<r}(1-q^{i-n})^2}{\prod_{i<r}(1-q^{i-r})}
=q^{-k^2}\,\frac{\prod_{j=k+1}^{n}(1-q^{-j})^2}{\prod_{j=1}^{n-k}(1-q^{-j})}.
\]
As $n\to\infty$ the products converge, because $\sum q^{-j}<\infty$. The limit is $\prod_{j>k}(1-q^{-j})^2/\prod_{j\ge1}(1-q^{-j})$, which is the stated limit.

(c) Set $c_q=\prod_{j\ge1}(1-q^{-j})>0$. The probabilities at stage $n$ are at most $q^{-k^2}/c_q$, and that bound is summable. They add to $1$ for each $n$. So the tail $\sum_{k>K}$ is at most $\frac1{c_q}\sum_{k>K}q^{-k^2}$ for every $n$. The limit may be taken term by term, and the limits add to $1$:
\[
\sum_{k\ge0}q^{-k^2}\frac{\prod_{j>k}(1-q^{-j})}{\prod_{j\le k}(1-q^{-j})}=1,\quad\text{that is}\quad
\sum_k\frac{x^{k^2}}{(x)_k^2}=\frac1{\prod_j(1-x^j)}\ \text{ at }x=\tfrac1q,
\]
after dividing by $c_q$ and writing $\prod_{j>k}(1-x^j)=c_q/(x)_k$. Both sides are power series in $x$ with non-negative integer coefficients, convergent for $|x|<1$. They agree at $x=\frac1q$ for every prime power $q$, and those points form a sequence tending to $0$. The identity theorem (Theorem~\ref{thm:identity}, on the unit disc) says the two series are equal.

\tools Volume~3, Chapter~9 (finite fields, and the count of bases of a vector space over $\mathbb F_q$); this chapter (the identity theorem).
\leads A partition has a largest square $k\times k$ in its Young diagram, the Durfee square. The two pieces outside that square are partitions into parts at most $k$. That is the combinatorial proof: $\frac{x^{k^2}}{(x)_k^2}$ counts the partitions whose Durfee square is $k$. For $q=2$ the limit law is $0.2888$, $0.5776$, $0.1284$, $0.0052$ for coranks $0,1,2,3$. The same shape appears, conjecturally, in the Cohen--Lenstra law for class groups of number fields.
\end{solution}

\begin{solution}{pb:dedekind}
\idea A Dedekind sum is the correlation of the sawtooth with a permuted copy of itself. In the Fourier basis of $\Z/k\Z$ the sawtooth becomes a sum of cotangents, so $s(h,k)$ is a cotangent sum, and cotangent sums are residues. The function $\cot(\pi z)\cot(\pi hz)\cot(\pi kz)$ has residues that produce $s(h,k)$, $s(k,h)$, and an explicit term at $z=1$. Its integral around a period rectangle can be computed. That computation is the reciprocity law, and the law makes $s(h,k)$ computable along Euclid's algorithm.

\writeup (a) Set $w=e^{2\pi ij/k}$, with $1\le j\le k-1$ and $0<r<k$. Then $\cot\frac{\pi j}k\sin\frac{2\pi jr}k=\frac{(w+1)(w^r-w^{-r})}{2(w-1)}=\frac12(w+1)\sum_{m=-r}^{r-1}w^m=\frac12\bigl(w^{-r}+w^r+2\sum_{|m|<r}w^m\bigr)$. Sum over $j$. The geometric sum $\sum_{j=1}^{k-1}w^m=k[k\mid m]-1$ gives $\frac12\bigl(-2+2(k-1)-2(2r-2)\bigr)=k-2r$, and $-\frac{k-2r}{2k}=\frac rk-\frac12=((\frac rk))$. For $r=0$ both sides vanish. Then use $\sum_{r\bmod k}\sin\frac{2\pi jr}k\sin\frac{2\pi j'r}k=\frac k2([j\equiv j']-[j\equiv-j'])$ and the substitution $j'\mapsto hj'$:
\[
s(h,k)=\frac1{4k^2}\sum_{j,j'}\cot\frac{\pi j}k\cot\frac{\pi j'}k\sum_r\sin\frac{2\pi jr}k\sin\frac{2\pi hj'r}k=\frac1{4k}\sum_{j=1}^{k-1}\cot\frac{\pi j}k\cot\frac{\pi hj}k,
\]
The two Kronecker terms contribute equally, because $\cot$ is odd.

(b) The function $\Phi(z)=\cot(\pi z)\cot(\pi hz)\cot(\pi kz)$ has period $1$, so the vertical sides cancel. As $Y\to\infty$, one has $\cot(\pi(x\pm iY))\to\mp i$ uniformly. The top side, traversed leftwards, tends to $-(-i)^3=-i$, and the bottom side tends to $i^3=-i$. The integral tends to $-2i$, and the residues add up to $-\frac1\pi$. At $z=\frac jh$ with $0<j<h$, only $\cot(\pi hz)$ has a pole, with residue $\frac1{\pi h}$. Those residues add up to $\frac1{\pi h}\sum_j\cot\frac{\pi j}h\cot\frac{\pi kj}h=\frac4\pi s(k,h)$, and the residues at $\frac jk$ add up to $\frac4\pi s(h,k)$. At $z=1$, the expansion $\cot(\pi u)=\frac1{\pi u}-\frac{\pi u}3+O(u^3)$ gives $\Phi(1+u)=\frac1{\pi^3hku^3}(1-\frac{\pi^2u^2}3)(1-\frac{\pi^2h^2u^2}3)(1-\frac{\pi^2k^2u^2}3)+\dots$, with residue $-\frac{1+h^2+k^2}{3\pi hk}$. So $\frac4\pi(s(h,k)+s(k,h))-\frac{1+h^2+k^2}{3\pi hk}=-\frac1\pi$, which is the law.

(c) Directly, $s(1,k)=\sum_{r=1}^{k-1}(\frac rk-\frac12)^2=\frac{(k-1)(2k-1)}{6k}-\frac{k-1}2+\frac{k-1}4=\frac{(k-1)(k-2)}{12k}$. Write $\{\frac{hr}k\}=\frac{hr\bmod k}k$ and use $\sum_{r=1}^{k-1}\{\frac{hr}k\}=\frac{k-1}2$. Then $s(h,k)=\frac1{k^2}\sum_rr\,(hr\bmod k)-\frac{k-1}4$. Since $r(hr\bmod k)\equiv hr^2\pmod k$, one has $6\sum_rr(hr\bmod k)\equiv h(k-1)k(2k-1)\equiv0\pmod k$, so $6k\,s(h,k)=\frac6k\sum_rr(hr\bmod k)-\frac{3k(k-1)}2$ is an integer. Finally Cauchy--Schwarz, and the fact that $r\mapsto hr$ permutes the residues, give $|s(h,k)|\le\sum_r((\frac rk))^2=s(1,k)$.

(d) The sum $s(h,k)$ depends on $h$ only modulo $k$. Let $0<h<k$ and write $k=a_1h+r_1$, so that $\frac{r_1}h=[0;a_2,\dots,a_n]$. Part (b) gives $|s(h,k)|\le|s(r_1,h)|+\frac1{12}(\frac kh+\frac hk+\frac1{hk})+\frac14\le|s(r_1,h)|+\frac{a_1}{12}+1$, because $\frac kh<a_1+1$, $\frac hk<1$ and $\frac1{hk}\le\frac12$. Iterate down to $s(0,1)=0$. Then $|s(h,k)|\le\frac1{12}\sum a_i+n$. The same recursion, kept with signs, computes $s(h,k)$ exactly, and Euclid's algorithm has $n=O(\log k)$ steps.

(e) If $\frac{s(h,k)}{\log k}$ were exactly Cauchy, then $P(|s|>t\log k)\approx\frac c t$, and the mean of $|s|/\log k$ would diverge logarithmically. The bound $|s(h,k)|\le s(1,k)<\frac k{12}$ cuts the tail at $t\approx\frac k{\log k}$, and $\mathbb E\frac{|s|}{\log k}\approx\int_1^{k/\log k}\frac c t\dd t\approx c\log k$. So the average of $|s|$ grows like $\log^2k$, as the numbers show: the ratios to $\log^2k$ are $0.045$, $0.043$, $0.044$, $0.044$. The large values come from an $h$ with one large partial quotient, as (d) suggests: $h=1$, $h=2$, or $h$ close to $\frac ka$.

\tools{Chapter~\ref{ch:fft} (discrete Fourier analysis on $\Z/k\Z$); this chapter (residues); Volume~2, Chapters~1 and~9 (Euclid's algorithm, continued fractions).}

\leads Dedekind met these sums in 1892, in the transformation law of $\log\eta(\tau)$ under $\tau\mapsto\frac{a\tau+b}{c\tau+d}$. That law involves $s(d,c)$. Theorem~\ref{thm:eta} of Chapter~\ref{ch:hr} is the case $c=1$, where $s(0,1)=0$. The proof of reciprocity given here is one of several in the monograph of Rademacher and Grosswald.
\end{solution}

\begin{solution}{pb:zigzagpole}
\idea The nearest singularity of a generating function sets the growth of its coefficients. Subtract that singularity. What remains is holomorphic on a larger disc, so its coefficients are smaller by a geometric factor. Here the nearest pole is at $\frac\pi2$, and the next one is at $-\frac{3\pi}2$, three times as far.

\writeup (a) The zeros of $\cos$ are the points $\frac\pi2+j\pi$, and each is simple. At $z=\frac\pi2$, one has $1+\sin z=2$ and $\cos'=-\sin=-1$, so $g$ has a simple pole with residue $-2$. The principal part is $\frac{-2}{z-\pi/2}=\frac4{\pi-2z}$. At $z=-\frac\pi2$, the factor $1+\sin z$ has a double zero (its derivative $\cos z$ vanishes there too) and $\cos$ has a simple zero, so the singularity is removable. Every other zero of $\cos$ has modulus at least $\frac{3\pi}2$. So $h(z)=g(z)-\frac4{\pi-2z}$ extends to a holomorphic function on $|z|<\frac{3\pi}2$.

(b) Cauchy's estimates (Corollary~\ref{cor:radius}), on a circle of radius $\rho<\frac{3\pi}2$, say the Taylor coefficients of $h$ at $0$ are $O(\rho^{-n})$. And $[z^n]\frac4{\pi-2z}=\frac4\pi\bigl(\frac2\pi\bigr)^n=2\bigl(\frac2\pi\bigr)^{n+1}$. So $\frac{A_n}{n!}=2\bigl(\frac2\pi\bigr)^{n+1}+O(\rho^{-n})$, with relative error $O\bigl((\frac\pi{2\rho})^n\bigr)$. The ratio $\frac\pi{2\rho}$ can be any number greater than $\frac13$.

(c) Compute $\frac{A_{10}}{10!}=\frac{50521}{3628800}=0.0139222$ and $2(\frac2\pi)^{11}=0.0139223$. The relative error is $-5.6\cdot10^{-6}$, and $-3^{-11}=-5.6\cdot10^{-6}$ as well.

\tools{Chapter~\ref{ch:generating} (the generating function of $A_n$); this chapter (poles and Cauchy's estimates).}
\leads Subtract the poles one after another and the exact formula appears: $\frac{A_n}{n!}=2(\frac2\pi)^{n+1}\sum_{k\ge0}(-1)^{k(n+1)}(2k+1)^{-(n+1)}$. The error is about $(-1)^{n+1}3^{-(n+1)}$, as (c) already shows. For odd $n$ the sum is $(1-2^{-n-1})\zeta(n+1)$. The tangent numbers $A_{2m-1}$ are the coefficients behind $\zeta(2m)$, which Chapter~\ref{ch:zeta} computes from Bernoulli numbers.
\end{solution}

\begin{solution}{ex:jensen}
\idea If $f$ had no zeros, $\log|f|$ would be the real part of a holomorphic function, and its mean over a circle would equal its value at the centre. Each zero breaks that. Divide the zero out by a factor of modulus $1$ on the circle, so the mean over the circle does not see the factor. At the centre, read what the division cost.

\writeup (a) The quotient $\frac{g'}g$ is holomorphic on the disc. The disc is star-shaped, so Theorem~\ref{thm:cauchy} supplies a primitive $L_0$. Then $(ge^{-L_0})'=g'e^{-L_0}-g\frac{g'}ge^{-L_0}=0$, so $ge^{-L_0}$ is a constant $c$. Integrate the derivative along segments from $0$ (Proposition~\ref{prop:pathint}(b)), and note $c\ne0$. Write $c=e^\lambda$ and $L=L_0+\lambda$. Then $g=e^L$, and $\log|g|=\operatorname{Re}L$. The mean value property (Theorem~\ref{thm:cauchyformula} at the centre) gives $L(0)=\frac1{2\pi}\int_0^{2\pi}L(re^{it})\dd t$. Take real parts.

(b) The factor $B_a$ is holomorphic for $|z|<\frac{R^2}{|a|}$, a disc larger than $\overline D(0,R)$, and its only zero is $a$, a simple one. For $|z|=R$, one has $R^2-\overline az=z\overline z-\overline az=z\,\overline{(z-a)}$, of modulus $R|z-a|$, so $|B_a(z)|=1$. And $|B_a(0)|=\frac{|a|}R$.

(c) The zeros of $f$ are isolated, so a slightly larger closed disc contains only finitely many of them. Choose $R'>R$ so that $D(0,R')$ lies inside the domain of $f$ and of every $B_{a_k}$, and so that $f$ has no zero with $R\le|z|<R'$. Set $g=f/\prod_kB_{a_k}$. At each $a_k$ the zeros of the numerator and of the denominator have the same order, so $g$ extends to a holomorphic function on $D(0,R')$ with no zeros (Theorem~\ref{thm:identity}(a) and Theorem~\ref{thm:removable}). Part (a) at $r=R$, together with (b), gives
\[
\frac1{2\pi}\int_0^{2\pi}\log\bigl|f(Re^{it})\bigr|\dd t=\frac1{2\pi}\int_0^{2\pi}\log\bigl|g(Re^{it})\bigr|\dd t=\log|g(0)|=\log|f(0)|+\sum_k\log\frac R{|a_k|}.
\]

(d) Choose $R>r$ with no zero of $f$ on $|z|=R$. All but countably many radii qualify, because the zeros are isolated and therefore countable. In Jensen's formula the left side is at most $\log C+AR$, one has $\log|f(0)|=0$, every term of the sum is non-negative, and each of the $N(r)$ zeros with $|a_k|\le r$ contributes at least $\log\frac Rr$. So $N(r)\log\frac Rr\le AR+\log C$. Both sides are continuous in $R$, so the inequality holds for every $R>r$. Take $R=er$.

For $f(z)=\frac{\sin\pi z}{\pi z}$ (with $f(0)=1$), the series gives $\bigl|\frac{\sin w}w\bigr|\le\sum_{k\ge0}\frac{|w|^{2k}}{(2k+1)!}\le\sum_{k\ge0}\frac{|w|^{2k}}{(2k)!}\le e^{|w|}$, so $|f(z)|\le e^{\pi|z|}$: here $C=1$ and $A=\pi$. The bound is $N(r)\le e\pi r\approx8.54\,r$. The true count is the number of non-zero integers of modulus at most $r$, which is $2\lfloor r\rfloor$. The bound has the right order. An entire function of exponential growth cannot have more than about $r$ zeros in a disc of radius $r$. In that sense the sine is as full of zeros as its growth allows, up to a constant.

\tools{This chapter: Cauchy's theorem on a disc, the mean value property, and the factorisation at a zero.}
\leads Jensen published the formula in 1899. It is the first step of Hadamard's factorisation of entire functions (Volume~8, Chapter~14). Applied to the zeta function on discs centred to the right of the critical strip, it bounds the number of zeros of $\zeta$ in a box of height $1$ at height $T$ by a constant times $\log T$. That is the first step toward the count ``about $\frac T{2\pi}\log T$ zeros up to height $T$'' of Riemann and von Mangoldt. This volume does not need that count.
\end{solution}

\begin{solution}{ex:rw09a}
\idea Proposition~\ref{prop:pathint}(c) already names the answers. The exercise asks you to see them from the parametrisation $z=e^{it}$.
\writeup Let $z=e^{it}$, so $\mathrm dz=ie^{it}\,\mathrm dt$ and $t$ runs from $0$ to $2\pi$. Then
\[
\int_{C(0,1)}z^2\dd z=\int_0^{2\pi}e^{2it}\cdot ie^{it}\dd t=i\int_0^{2\pi}e^{3it}\dd t=\Bigl[\frac{e^{3it}}3\Bigr]_0^{2\pi}=0,
\]
because $e^{6\pi i}=1=e^{0}$. For the second integral,
\[
\int_{C(0,1)}\frac{\mathrm dz}z=\int_0^{2\pi}i\dd t=2\pi i.
\]
The real number $2\pi$ is $6.283185\ldots$. The same computation on the circle of radius $2$ still gives $2\pi i$ for $\int\mathrm dz/z$, since the $r$ in the numerator and the denominator cancel.
\end{solution}

\begin{solution}{ex:rw09b}
\idea The unit circle is the positive boundary of a disc. Write the integral as a residue at the one root of a quadratic that lies inside that disc.
\writeup Set $z=e^{i\theta}$, so $\mathrm d\theta=\mathrm dz/(iz)$ and $\cos\theta=(z+1/z)/2$. Then
\[
2+\cos\theta=\frac{z^2+4z+1}{2z},\qquad \frac{\mathrm d\theta}{2+\cos\theta}=\frac{2\,\mathrm dz}{i(z^2+4z+1)}.
\]
The roots of $z^2+4z+1$ are $-2\pm\sqrt3$. The root inside the unit circle is $z_0=-2+\sqrt3$, because $|-2+\sqrt3|=2-\sqrt3\approx0.2679<1$ and the product of the roots is $1$. The residue of $2\big/\bigl(i(z^2+4z+1)\bigr)$ at $z_0$ is
\[
\frac{2}{i(2z_0+4)}=\frac{1}{i(z_0+2)}=\frac{1}{i\sqrt3}.
\]
Multiply by $2\pi i$: the integral equals $2\pi/\sqrt3$. Numerically, $2\pi/\sqrt3=3.6275987\ldots$. The eight-point sum is
\[
\frac{2\pi}8\sum_{k=0}^{7}\frac1{2+\cos(2\pi k/8)}=3.6277915\ldots,
\]
which differs from $2\pi/\sqrt3$ by about $1.93\times10^{-4}$.
\end{solution}

\begin{solution}{ex:rw09c}
\idea The denominator $(1+z^2)^2=(z-i)^2(z+i)^2$ has a pole of order $2$ at $i$. Proposition~\ref{prop:rescalc}(a) with $m=2$ is one differentiation.
\writeup The only pole in the upper half-plane is $i$. For a pole of order $2$,
\[
\operatorname{Res}_i\frac1{(1+z^2)^2}=\lim_{z\to i}\frac{\mathrm d}{\mathrm dz}(z+i)^{-2}=-2(z+i)^{-3}\Big|_{z=i}=-2(2i)^{-3}.
\]
Now $(2i)^3=8i^3=-8i$, so the residue is $-2/(-8i)=1/(4i)=-i/4$. The integral equals $2\pi i\cdot(-i/4)=\pi/2$. To eight decimals, $\pi/2=1.57079633$.

The same number comes from the real antiderivative. Differentiate $\frac12\bigl(x/(1+x^2)+\arctan x\bigr)$ and you recover $1/(1+x^2)^2$. From $-\infty$ to $\infty$ the boundary term $\frac x{1+x^2}$ vanishes, and $\arctan$ contributes $\frac12\bigl(\frac\pi2-(-\frac\pi2)\bigr)=\pi/2$.
\end{solution}

\begin{solution}{ex:rw09d}
\idea Theorem~\ref{thm:cauchy} builds the primitive by integrating along the segment from the star centre. Here the centre is $1$ and the function is $1/w$, so the integral should be $\operatorname{Log} i-\operatorname{Log} 1$.
\writeup Parametrise $[1,i]$ by $w(t)=(1-t)+it$, $t\in[0,1]$, so $w'(t)=-1+i$. The integral is
\[
\int_0^1\frac{-1+i}{(1-t)+it}\dd t.
\]
The denominator has modulus squared $(1-t)^2+t^2=2t^2-2t+1$. Multiply numerator and denominator by the conjugate $1-t-it$:
\[
(-1+i)(1-t-it)=-1+2t+i.
\]
The real part of the integrand is $(2t-1)/(2t^2-2t+1)$, whose antiderivative is $\frac12\log(2t^2-2t+1)$. From $0$ to $1$ that logarithm returns to $\log1$ at both ends, so the real part integrates to $0$. The imaginary part is $1/(2t^2-2t+1)$. Write $2t^2-2t+1=2\bigl(t-\frac12\bigr)^2+\frac12$. The substitution $v=2t-1$ gives
\[
\int_0^1\frac{\mathrm dt}{2t^2-2t+1}=\int_{-1}^{1}\frac{\mathrm dv}{v^2+1}=\arctan1-\arctan(-1)=\frac\pi2.
\]
The integral equals $i\pi/2$, which is $\operatorname{Log} i$. Numerically $i\pi/2\approx1.570796i$. The midpoint rule with panels at $t=\frac18,\frac38,\frac58,\frac78$ returns $1.581176i$. The two imaginary numbers differ by about $0.0104$.
\end{solution}

\section*{Sources and further reading}

\begin{itemize}
\item E.~Stein and R.~Shakarchi, \emph{Complex Analysis}, Princeton University Press, 2003, Chapters~1 to~3 (Goursat,
Cauchy, residues, Rouch\'e) and~6 (the Gamma function).
\item L.~Ahlfors, \emph{Complex Analysis}, 3rd ed., McGraw-Hill, 1979. The classic, with the homology form of Cauchy's theorem.
\item T.~Needham, \emph{Visual Complex Analysis}, Oxford University Press, 1997. Pictures for everything in this chapter,
including winding numbers and Rouch\'e's ``dog on a leash''.
\item E.~Goursat, \emph{Sur la d\'efinition g\'en\'erale des fonctions analytiques, d'apr\`es Cauchy}, Transactions of the AMS 1
(1900), 14--16. Theorem~\ref{thm:goursat} without continuity of $f'$.
\item A.-L.~Cauchy, \emph{M\'emoire sur les int\'egrales d\'efinies prises entre des limites imaginaires}, Paris, 1825. The birth
of contour integration.
\item R.~Remmert, \emph{Wielandt's theorem about the $\Gamma$-function}, American Mathematical Monthly 103 (1996), 214--220.
Exercise~\ref{ex:wielandt}.
\item The epigraph: P.~Painlev\'e, \emph{Analyse des travaux scientifiques}, Gauthier-Villars, 1900; J.~Hadamard repeated the
thought in \emph{An Essay on the Psychology of Invention in the Mathematical Field}, Princeton University Press, 1945.
\item H.~Rademacher and E.~Grosswald, \emph{Dedekind Sums}, Carus Mathematical Monographs 16, Mathematical Association of
America, 1972; I.~Vardi, \emph{Dedekind sums have a limiting distribution}, International Mathematics Research Notices 1993,
1--12, for Exercise~\ref{pb:dedekind}.
\item J.~L.~W.~V.~Jensen, \emph{Sur un nouvel et important th\'eor\`eme de la th\'eorie des fonctions}, Acta Mathematica 22 (1899),
359--364, for Exercise~\ref{ex:jensen}.
\item D.~S.~Mitrinovi\'c and J.~D.~Ke\v cki\'c, \emph{The Cauchy Method of Residues}, Reidel, 1984. Hundreds of integrals and
sums computed by the methods of this chapter.
\end{itemize}
